\documentclass[11pt]{amsart}
\usepackage[margin=1.1in]{geometry}
\usepackage{amsmath,amssymb,amsthm,mathtools}
\usepackage{booktabs,array,longtable}
\usepackage[T1]{fontenc}
\usepackage{lmodern}
\usepackage{url}
\usepackage[hidelinks]{hyperref}

\newtheorem{theorem}{Theorem}[section]
\newtheorem{proposition}[theorem]{Proposition}
\newtheorem{lemma}[theorem]{Lemma}
\newtheorem{corollary}[theorem]{Corollary}
\newtheorem{conjecture}[theorem]{Conjecture}
\theoremstyle{definition}
\newtheorem{definition}[theorem]{Definition}
\newtheorem{hypothesis}[theorem]{Hypothesis}
\newtheorem{example}[theorem]{Example}
\newtheorem{remark}[theorem]{Remark}
\numberwithin{equation}{section}
\raggedbottom
\emergencystretch=3em

\newcommand{\re}{\operatorname{Re}}
\newcommand{\im}{\operatorname{Im}}
\newcommand{\ZZ}{\mathbb Z}
\newcommand{\QQ}{\mathbb Q}
\newcommand{\RR}{\mathbb R}
\newcommand{\CC}{\mathbb C}
\newcommand{\NN}{\mathbb N}
\newcommand{\TT}{\mathbb T}
\newcommand{\cA}{\mathcal A}
\newcommand{\cC}{\mathcal C}
\newcommand{\hA}{\lambda_{\cA}}
\newcommand{\dlow}{\underline d}
\newcommand{\dup}{\overline d}
\newcommand{\eps}{\varepsilon}
\newcommand{\spanQ}[1]{\langle#1\rangle_{\QQ}}
\newcommand{\file}[1]{\path{#1}}
\newcommand{\repo}{\texttt{2026-11-number-theory-q-series/04-chebyshev-bias-grh-only}}
\newcommand{\rd}[1]{\textup{[#1]}}

\begin{document}

\title[Chebyshev's bias without linear independence: the atom at zero]{Chebyshev's bias without linear independence:\\ certified bounds, the possible atom at zero, and barriers}

\author{Raghav Sharma}
\email{raghavsharma1729@gmail.com}
\date{27 November 2026}
\subjclass[2020]{Primary 11N13; Secondary 11M26, 42A75, 11Y35, 60E10}
\keywords{Chebyshev's bias, prime number races, logarithmic density, generalized Riemann hypothesis, linear independence hypothesis, almost periodic functions, limiting distribution, atoms, integer multiplicities, ball arithmetic}

\begin{abstract}
For $\chi=\chi_{-4}$ it is not known whether, under the generalized Riemann hypothesis (GRH) for $L(s,\chi)$ alone, the set $\{x:\pi(x;4,3)>\pi(x;4,1)\}$ has a logarithmic density: without a linear-independence hypothesis on the zeros, the limiting distribution might have an atom at $0$. We do not prove that the density exists. Assuming only GRH for the character concerned, we prove: the lower logarithmic density of this set exceeds $0.910924861$, and that of $\{x:\pi(x;3,2)>\pi(x;3,1)\}$ exceeds $0.934761398$ (certified in ball arithmetic); a time-reversal identity on the closure torus, which makes the limiting law symmetric about its mean up to an explicit error; that the mass of the possible atom at $0$ equals the jump of the bias as a function of a weight $p^m$, so that the density exists for all but countably many weights; and a reduction of the problem to a uniform small-ball bound for finite trigonometric models with integer multiplicities. We prove existence of the density under several hypotheses on the zeros that are weaker than linear independence, and under a prime-side band-decoupling hypothesis. Finally we construct models that share the structure used by these arguments, including positive integer multiplicities obeying a zero-counting law with a large constant, whose laws do have an atom at $0$; there the counting constant is astronomically large ($\gtrsim10^{14}$), and the true constant and the case of bounded multiplicities remain open.
\end{abstract}

\maketitle

\noindent\fbox{\parbox{\dimexpr\linewidth-2\fboxsep-2\fboxrule}{\small
\textbf{Disclosure.} This paper was produced by AI (Anthropic's Claude models), directed by the author. Its proofs and computations were checked by independent AI referees, each instructed to find errors (Section~\ref{sec:review}); every result stated here passed at least one such read, and the number of reads is recorded for each. The paper has \emph{not yet been refereed by a human expert}. The author takes responsibility for the content and welcomes corrections at raghavsharma1729@gmail.com. Code, certificates and referee reports: \url{https://github.com/raghavs1729/Research-Math}, directories \repo{} and \texttt{2026-11-number-theory-q-series/05-weighted-races-grh} (certified zero lists and variance and non-resonance certificates).}}

\medskip
\tableofcontents

\section{Introduction}\label{sec:intro}

\subsection{The problem}
Chebyshev observed in 1853 that primes $\equiv3\pmod4$ seem to outnumber primes $\equiv1\pmod4$. Rubinstein and Sarnak \cite{RS} made this precise: under the generalized Riemann hypothesis (GRH) for $L(s,\chi_{-4})$ and the linear independence hypothesis (LI: the positive ordinates of the zeros are linearly independent over $\QQ$), the set $\{x:\pi(x;4,3)>\pi(x;4,1)\}$ has logarithmic density $0.9959\ldots$. GRH alone already gives a limiting logarithmic distribution $\nu$ of the normalised race $(\log x/\sqrt x)(\pi(x;4,3)-\pi(x;4,1))$ \cite{RS}, which is the push-forward of Haar measure on the closure of a one-parameter subgroup of an infinite-dimensional torus (\cite{ANS,Devin,Bailleul}; \cite[Theorem~1.3]{GO}). The density exists if $\nu(\{0\})=0$, and LI is used to exclude an atom at $0$. Rubinstein and Sarnak note (\cite[p.~174]{RS}) that ``assuming only GRH we don't know that $\delta(P_{q;a_1,\dots,a_r})$ exists''. As far as we know this is still open. Several weakenings of LI suffice (Devin \cite[Theorems~2.2, 5.4]{Devin}, Martin--Ng \cite{MN}, Bailleul \cite{Bailleul}); none of them is known to hold. Item~1 of the problem list of Hamieh, Kadiri, Martin and Ng \cite{HKMN} (posed by Martin and Silberman) asks to prove that $\zeta(s)$ has at least one irrational ordinate, and that the positive ordinates of $\zeta(s)$ are not all rational multiples of one another; the analogous statements for $L(s,\chi)$ are equally open. Problem~13 of the same list asks to prove unconditionally that ties $\pi(x;q,a)=\pi(x;q,b)$ have density $0$; our results bear on this only conditionally on GRH (Corollary~\ref{cor:ties}). Under GRH and weakened forms of LI, density $0$ for ties follows from the continuity results of Devin \cite[Theorem~2.2]{Devin}, \cite[Corollary~2]{Devin2}.

\textbf{We do not prove that the logarithmic density of $\{x:\pi(x;4,3)>\pi(x;4,1)\}$ exists under GRH alone.} This paper collects what we could prove in that direction, and why the methods available to us cannot finish the job.

\subsection{Setting}
Let $\chi$ be a real primitive character modulo $q$, $\chi^*=-\chi$, and for real $m>-\frac12$ put $M=m+\frac12$ and
\[
D_m(x)=\sum_{p\le x}\chi^*(p)p^m,\qquad E_M(x)=2M\,x^{-M}(\log x)\,D_m(x).
\]
For $q=4$, $m=0$: $D_0(x)=\pi(x;4,3)-\pi(x;4,1)$ and $E_{1/2}(x)=(\log x/\sqrt x)D_0(x)$; for $\chi=\chi_{-3}$, $D_0(x)=\pi(x;3,2)-\pi(x;3,1)$. The weight $p^m$ interpolates between the classical race ($m=0$) and the race weighted by $p^{-1/2}$ ($m\to-\frac12$), for which Hayani \cite[Corollary~1.2]{Hayani} proved, under GRH alone, that the lead set has natural density $1$ ($q=4$); Sheth \cite{Sheth} has closely related results under GRH, and the $p^{-1/2}$-weighted framework goes back to Aoki and Koyama \cite{AK}. GRH$(\chi)$ means that all nontrivial zeros of $L(s,\chi)$ have real part $\frac12$ (no hypothesis on any other $L$-function, none on $\zeta$); NCZ$(\chi)$ means $L(\frac12,\chi)\ne0$, which is a theorem for $q=3,4$. Under GRH and NCZ, $E_M(e^y)$ has a limiting distribution $\nu_M$ in $y$ (Section~\ref{sec:setting}). Results conditional on GRH say so in their statements. A Lean proof of GRH for Dirichlet $L$-functions was announced on 2026-11-20; it is unrefereed and is not used anywhere in this paper. It would in any case not touch the question studied here, which concerns linear relations among the zeros.

\subsection{Results}
All statements below assume GRH$(\chi)$ and NCZ$(\chi)$ unless stated otherwise, and none assumes LI.

\emph{Certified bounds} (Theorem~\ref{thm:cert}). Under GRH$(\chi_{-4})$ alone, the lower logarithmic density of $\{x:\pi(x;4,3)>\pi(x;4,1)\}$ is $>0.910924861$; under GRH$(\chi_{-3})$ alone, that of $\{x:\pi(x;3,2)>\pi(x;3,1)\}$ is $>0.934761398$. The previous certified values without LI were $0.884276434$ and $0.911240824$ \cite{GO}. So $\nu_{1/2}(\{0\})<0.0891$, resp.\ $<0.0653$.

\emph{Time reversal} (Theorem~\ref{thm:reversal}). On the closure torus $\cA$, with $F_M$ the function whose law is $\nu_M$, one has $F_M(\theta)+F_M(\theta^{-1})=2+R(\theta)$ almost everywhere, with $|R|\le\sigma_M^2=\sum_{\gamma>0}\frac{8M^2}{M^2+\gamma^2}$ (sum over the positive ordinates with multiplicity; a closed form). So $\nu_M$ is symmetric about its mean $1$ up to $\sigma_M^2$ in the $L^\infty$-Wasserstein distance, without LI, and $\nu_M((-\infty,0])\le\frac12$ whenever $\sigma_M^2<2$.

\emph{Atom $=$ jump} (Theorem~\ref{thm:jump}). For $M$ in the interior of an explicit interval containing $\frac12$, $\nu_M(\{0\})$ equals the jump at $m$ of the nonincreasing function $m\mapsto\dlow\{x:D_m(x)>0\}$. Hence the logarithmic density of $\{D_m>0\}$ exists for all but countably many $m\in(-\frac12,1.3239]$ ($q=4$), and the classical race has a density as soon as the bias is continuous in the weight at $m=0$.

\emph{Sufficient conditions} (Section~\ref{sec:sufficient}). The density exists for every $m$ under GRH together with any one of the following hypotheses on the ordinates, each implied by LI: TS (the ordinates above some height do not span all ordinates over $\QQ$); FD (finite-level directions); TS$_\ZZ$ (the ordinates above height $T$ generate subgroups of unbounded index); V (some ordinate never carries half of the $\ell^1$-mass of an integer relation); GLI (graded linear independence: relations reaching height $t$ have length $\ge\psi(t)$ with $\psi(t)\log t/t\to\infty$); or a prime-side band-decoupling hypothesis on mixed moments of windowed prime sums (implied by LI together with simplicity of the zeros). As conditions on sets of frequencies, TS and V are strictly weaker than the hypothesis of Devin's Theorem~2.2, and FD is strictly weaker than each hypothesis of this kind that we found in the literature. Conversely, an atom forces strong arithmetic structure: every ordinate has a multiple $n\gamma$, $n\le1/\nu_M(\{c\})$, that is an integer combination of ordinates above any given height, and integer relations of length $\le Ct/\log t$ reaching height $t$ occur infinitely often.

\emph{Reduction} (Lemma~\ref{lem:reduction}, Conjecture~\ref{conj:U}). The density at $m=0$ follows from a uniform small-ball bound for finite trigonometric polynomials with positive integer multiplicities obeying the zero-counting law, on arbitrary subtori; in this form the unknown group generated by the zeros disappears. We prove this bound on the subclass of models with a low Newton vertex (Corollary~\ref{cor:Ucase}).

\emph{Barriers} (Section~\ref{sec:barriers}). There are models with all the soft structure used above (a compact connected hull with dense flow, positive weights, the identities between weights, time reversal, $\sigma^2<2$ at the classical weight) whose law has an atom at $0$ (Proposition~\ref{prop:model73}, certified). Integer multiplicities forbid flat arcs (Theorem~\ref{thm:arcs}) and forbid atoms in structured classes (Theorem~\ref{thm:nested} and Section~\ref{sec:integrality}), but not in general: there are positive integer multiplicities $m_n\ll(\log n)^u$, for any $u>0$, whose model vanishes on a set of positive measure, and at the classical weight such models can have $\sigma^2<\frac12$ and obey the local zero-counting law, with a constant of size at least about $10^{14}$ (larger for the models with $\sigma^2<\frac12$) (Theorem~\ref{thm:noise}, Corollary~\ref{cor:noise}). The fixed-constant form (Conjecture~\ref{conj:U}) and the case of bounded multiplicities are untouched by this.

\subsection{Relation to the literature}
That a race against a \emph{level} has a density for all but countably many levels is classical; Devin \cite{Devin} attributes a statement of this type under RH alone to Kaczorowski and Ramar\'e \cite[Theorem~3]{KR} (we have not read it). What is specific to Theorem~\ref{thm:jump} is that the level is fixed at $0$ and the weight varies; it extends \cite[Theorem~7.1]{GO}, which gave only that ties are immaterial outside a countable set of weights. The mechanism of Theorems~\ref{thm:TS}, \ref{thm:FD} (move along a one-parameter subgroup of the hull on which the function is a non-constant trigonometric polynomial) is Devin's \cite{Devin}; what is added is that conditioning on the complementary coordinates replaces independence. The nearest prior result to Theorem~\ref{thm:reversal} that we found is Devin's exact symmetry \cite[Theorem~2.3]{Devin} under a parity hypothesis on the integer relations among the ordinates (also unproved). Chebyshev-type lower bounds without LI, with the multiplicity-squared variance, are Devin \cite[Corollary~5.8]{Devin} and Fiorilli \cite[Lemma~2.7]{Fiorilli}. In the papers we searched (listed in \file{density/LITERATURE.md}) we did not find the time-reversal identity, an explicit GRH-only lower bound for the bias of the $q=4$ race other than the author's \cite{GO}, or hypotheses of the TS, FD, V type; this is a statement about those searches only. No search was made for finite-index conditions or graded independence conditions, and no novelty is claimed for TS$_\ZZ$ and GLI. Lemma~\ref{lem:JN} is classical.

\subsection{Status of the results}
Every result stated as a theorem, proposition, lemma or corollary passed at least one adversarial read by an independent AI referee instructed to find errors (Section~\ref{sec:review}); the number of reads is given in brackets in each statement, e.g.\ \rd{2 reads}. Statements labelled \emph{Remark} or \emph{Heuristic} are not claimed as proved, except where a remark explicitly derives a consequence of a stated result. Many proofs are given in full; many others (in particular in Sections~\ref{sec:sufficient}--\ref{sec:barriers}) are outlined or only referenced, with complete proofs in the working notes cited at each place. File names refer to the directory \repo{} of \cite{Repo} unless they start with \texttt{05-}.

Table~\ref{tab:summary} lists the main results.

\begin{table}[ht]\small
\caption{Main results. ``GRH'' means GRH$(\chi)$ and NCZ$(\chi)$ for the character concerned (NCZ is a theorem for $q=3,4$).}\label{tab:summary}
\begin{tabular}{p{0.37\linewidth}p{0.2\linewidth}p{0.12\linewidth}p{0.17\linewidth}}\toprule
Statement&Hypotheses&Where&Reads\\\midrule
certified lower densities $0.910924861$ ($q=4$), $0.934761398$ ($q=3$)&GRH&Thm~\ref{thm:cert}&2, certified\\
time reversal; $\nu_M((-\infty,0])\le\frac12$ if $\sigma_M^2<2$&GRH&Thm~\ref{thm:reversal}&2\\
atom at $0$ $=$ jump in the weight; density for all but countably many weights&GRH&Thm~\ref{thm:jump}&2\\
density exists under TS, FD, tame directions&GRH $+$ hypothesis&\S\ref{sec:TS}&1\\
density exists under TS$_\ZZ$&GRH $+$ TS$_\ZZ$&Thm~\ref{thm:TSZ}&2\\
density exists under V; low vertices&GRH $+$ V / LV&Thm~\ref{thm:V}&1\\
density exists under GLI; what an atom forces&GRH $+$ GLI&Thm~\ref{thm:GLI}&2\\
prime-side duality; band decoupling&GRH $+$ prime-side&Thms~\ref{thm:P1}, \ref{thm:P7}&2\\
reduction to finite integer models; U on a subclass&GRH / none&\S\ref{sec:reduction}&1\\
positive-weight models with an atom, $M=\frac12$, $\sigma^2<0.55$&none&Prop~\ref{prop:model73}&1, certified\\
integrality: parity on nested combs, Toeplitz classes, arcs&none&\S\ref{sec:integrality}&1--2\\
integer models with an atom, $m_n\ll(\log n)^u$, counting constant $\gtrsim10^{14}$&none&Thm~\ref{thm:noise}&1\\\bottomrule
\end{tabular}
\end{table}

\section{Setting and quoted inputs}\label{sec:setting}

\subsection{The hull}
Let $\Gamma$ be the set of \emph{distinct} positive ordinates $\gamma$ of zeros $\frac12+i\gamma$ of $L(s,\chi)$, and $m_\gamma\ge1$ the multiplicity of $\gamma$. Since $\chi$ is real, the zeros are symmetric about the real axis. For $M>0$ put
\[
r_\gamma(M)=\frac{4M}{M+i\gamma},\qquad w_\gamma=\frac{8M^2}{M^2+\gamma^2}=2\re r_\gamma,\qquad s_\gamma=\frac{4M\gamma}{M^2+\gamma^2}=-\im r_\gamma,
\]
\[
\sigma_M^2=\sum_{\gamma\in\Gamma}m_\gamma w_\gamma,\qquad V_M=\sum_{\gamma\in\Gamma}m_\gamma^2w_\gamma=\tfrac12\sum_{\gamma\in\Gamma}m_\gamma^2|r_\gamma|^2,
\]
\[
\lambda_\chi(M,M')=\frac1{MM'}-\sum_{\gamma>0}\frac4{|M+i\gamma||M'+i\gamma|},
\]
the last sum with multiplicity. Under GRH and NCZ, $\sigma_M^2=4M(\Lambda'/\Lambda)(M+\frac12)$ with $\Lambda(s,\chi)$ the completed $L$-function \cite[\S1.1, (Q2)]{GO}; $V_M\ge\sigma_M^2$ with equality iff all zeros are simple. Let $\iota(y)=(e^{i\gamma y})_{\gamma\in\Gamma}\in\TT^\Gamma$, let $\cA$ be the closure of $\iota(\RR)$ (a compact connected metrizable abelian group) with Haar probability measure $\hA$, $e_\gamma$ the coordinate characters, $e_{-\gamma}=\bar e_\gamma$, and $\tau_s\theta=\theta\cdot\iota(s)$ the translation flow. Let $G\subset\RR$ be the subgroup generated by $\Gamma$; for $g=\sum n_\gamma\gamma\in G$ put $e_g=\prod e_\gamma^{n_\gamma}$. Define
\[
F_M=1+\re\sum_{\gamma\in\Gamma}m_\gamma r_\gamma(M)e_\gamma\ \in L^2(\cA),\qquad \|F_M-1\|_2^2=V_M,
\]
and let $\nu_M$ be the law of $F_M$ under $\hA$. For $T>0$ let $\Gamma_T=\Gamma\cap(0,T]$, $\Gamma_{>T}=\Gamma\setminus\Gamma_T$, $F_{M,T}=1+\re\sum_{\gamma\in\Gamma_T}m_\gamma r_\gamma e_\gamma$ and $V_M^{>T}=\frac12\sum_{\gamma>T}m_\gamma^2|r_\gamma|^2$. Let $M_\sigma(\chi)$ be defined by $\sigma^2_{M_\sigma}=2$; then $1.82391477\le M_\sigma(\chi_{-4})<1.82391478$ and $2.12775095\le M_\sigma(\chi_{-3})<2.12775096$ (certified, \cite{GO}).

For a measurable $B\subset[2,\infty)$ the lower and upper logarithmic densities are $\dlow(B)=\liminf_X(\log X)^{-1}\int_{[2,X]\cap B}dt/t$ and $\dup(B)$ with $\limsup$.

\subsection{Quoted inputs}
\begin{itemize}
\item[(I1)] (\cite[Lemma~1.2, Theorem~1.3]{GO}; cf.\ \cite{RS,ANS,Devin,Bailleul}.) Assume GRH$(\chi)$ and NCZ$(\chi)$. Time averages of continuous functions of $\iota(y)$ converge to Haar integrals; the characters $e_g$, $g\in G$, are orthonormal; the flow $\tau$ is ergodic; nonempty open sets have positive measure; and $E_M(e^y)$, as well as $E^\theta_M(e^y)=2Me^{-My}\sum_{p\le e^y}\chi^*(p)p^m\log p$, has limiting distribution $\nu_M$: $\frac1Y\int_0^Y\varphi(E_M(e^y))\,dy\to\int\varphi\,d\nu_M$ for bounded continuous $\varphi$. No LI is used: a character $\prod\theta_\gamma^{n_\gamma}$ is trivial on $\cA$ iff $\sum n_\gamma\gamma=0$. The proof rests on the explicit formula with error term from \cite{WR}.
\item[(I2)] (Pontryagin duality; \cite[2.1.2, 2.1.3]{Rudin}, numbering recalled.) $\cA=\{\theta\in\TT^\Gamma:\prod\theta_\gamma^{n_\gamma}=1\text{ whenever }\sum n_\gamma\gamma=0\}$, and its dual group is $\{e_g:g\in G\}\cong G$.
\item[(I3)] (Portmanteau.) If $\nu_Y\to\nu$ weakly, then $\liminf\nu_Y(U)\ge\nu(U)$ for open $U$ and $\limsup\nu_Y(C)\le\nu(C)$ for closed $C$. Consequently, for every $c$ and $\eps>0$,
\begin{equation}\label{eq:portmanteau}
\nu_M((c-\eps,c+\eps))\le\dlow\{|E_M-c|<\eps\}\le\dup\{|E_M-c|\le\eps\}\le\nu_M([c-\eps,c+\eps]).
\end{equation}
\item[(I4)] (\cite[Theorem~7.1]{GO}, \cite{MC}.) $\lambda_\chi(M,M')>0$ whenever $\sigma_M\sigma_{M'}<2$, in particular for $0<M<M'\le M_\sigma(\chi)$.
\item[(I5)] (Weil's explicit formula \cite[Theorem~5.12]{IK} (numbering recalled), for $\chi$ real primitive and a test function supported in $(0,\infty)$.) For $\eta\in C_c^\infty((0,\infty))$, $\sum_{\gamma\in\pm\Gamma}m_\gamma\int\eta(t)e^{i\gamma t}dt=-\sum_{n\ge2}\Lambda(n)\chi(n)n^{-1/2}\eta(\log n)+\int\eta(t)\Omega(t)\,dt$ with $\Omega\in C^\infty((0,\infty))$ coming from the Gamma factor.
\item[(I6)] (Tur\'an--Nazarov \cite{Nazarov}.) There is an absolute $A_N\ge1$ such that for every exponential polynomial $p(t)=\sum_{k=1}^nc_ke^{i\lambda_kt}$, every interval $I$ and every measurable $E\subset I$ of positive measure, $\sup_I|p|\le(A_N|I|/|E|)^{n-1}\sup_E|p|$. (Statement recalled; the numbering in the source was not re-read.)
\item[(I7)] (Zero counting, \cite[Ch.~16]{Davenport}.) $N(t+1)-N(t)\le A_0\log(t+3)$ for $t\ge0$, where $N(t)$ counts zeros with $0<\gamma\le t$ with multiplicity; $N(t)\sim\frac t{2\pi}\log t$. Hence $m_\gamma\le A_0\log(\gamma+3)$ and $V_M^{>T}\le C_1M^2\log^2(T+3)/T$.
\end{itemize}

Under GRH and NCZ, (I1) and (I3) give for every $m>-\frac12$, with $d(M):=\nu_M((0,\infty))$ and $\bar d(M):=\nu_M([0,\infty))$,
\begin{equation}\label{eq:sandwich}
d(M)\le\dlow\{D_m>0\}\le\dup\{D_m>0\}\le\dup\{D_m\ge0\}\le\bar d(M),\qquad d(M)\le\dlow\{D_m\ge0\}\le\bar d(M).
\end{equation}
In particular, if $\nu_M(\{0\})=0$ then the logarithmic densities of $\{D_m>0\}$ and $\{D_m\ge0\}$ exist and equal $d(M)$.

\begin{lemma}[\rd{1 read}]\label{lem:trigpoly}
A trigonometric polynomial $P=\sum_{g\in S}c_ge_g$ on $\cA$ ($S\subset G$ finite, not all $c_g=0$) vanishes only on a $\hA$-null set.
\end{lemma}

\begin{proof}
The subgroup $G_S$ generated by $S$ is free of finite rank $n$; let $g_1,\dots,g_n$ be a basis and $\pi=(e_{g_1},\dots,e_{g_n}):\cA\to\TT^n$. Its image is compact and connected and no nontrivial character of $\TT^n$ is identically $1$ on it, so $\pi$ is onto and $\pi_*\hA$ is Haar measure. Now $P=Q\circ\pi$ with $Q$ a nonzero trigonometric polynomial on $\TT^n$, whose zero set is Lebesgue-null (apply real analyticity to $|Q|^2$).
\end{proof}

\begin{lemma}[\rd{1 read}]\label{lem:analytic}
If $h$ is real-analytic and non-constant on an open interval $J$, the image of Lebesgue measure on $J$ under $h$ is absolutely continuous.
\end{lemma}

\begin{proof}
The zeros of $h'$ are isolated; between them $h$ is a diffeomorphism onto its image.
\end{proof}

\section{The atom at \texorpdfstring{$0$}{0} is the jump of the bias in the weight}\label{sec:jump}

Throughout this section assume GRH$(\chi)$ and NCZ$(\chi)$, and let $I$ be an interval of weights $M$ on which $\lambda_\chi(M,M')>0$ for all $M<M'$ in $I$; by (I4) one may take $I=(0,M_\sigma(\chi)]$, i.e.\ $-\frac12<m\le M_\sigma(\chi)-\frac12$ ($M_\sigma-\frac12=1.323914\ldots$ for $q=4$ and $1.627750\ldots$ for $q=3$), which contains the classical weight $m=0$ in its interior. For $0<M<M'$ put $\kappa=M'-M$ and
\[
G_{M,M'}(\theta)=1+\re\sum_{\gamma\in\Gamma}m_\gamma\frac{4MM'\theta_\gamma}{(M+i\gamma)(M'+i\gamma)},
\]
a continuous function on $\cA$. By the triangle inequality and (I4),
\begin{equation}\label{eq:Gbounds}
MM'\lambda_\chi(M,M')\le G_{M,M'}(\theta)\le2-MM'\lambda_\chi(M,M')\qquad(\theta\in\cA),
\end{equation}
and comparing Fourier coefficients ($\frac4{M+i\gamma}-\frac4{M'+i\gamma}=\frac{4\kappa}{(M+i\gamma)(M'+i\gamma)}$, constant terms $\frac1M-\frac1{M'}=\frac\kappa{MM'}$),
\begin{equation}\label{eq:transversal}
\frac{F_M}M-\frac{F_{M'}}{M'}=\frac\kappa{MM'}\,G_{M,M'}\qquad\hA\text{-a.e.}
\end{equation}
This is the hull form of the lead-inheritance theorem of \cite{MC}: for $m<m'$ in the range, $D_{m'}(x)\ge0$ implies $D_m(x)>0$ for large $x$.

\begin{theorem}[GRH, NCZ; \rd{2 reads}]\label{thm:jump}
\textup{(a)} For $M<M'$ in $I$: $\bar d(M')\le d(M)$. Hence $d$ and $\bar d$ are nonincreasing on $I$.

\textup{(b)} For every interior point $M_0$ of $I$: $d(M_0)=d(M_0^+)=\bar d(M_0^+)$, $\bar d(M_0)=\bar d(M_0^-)=d(M_0^-)$, and
\[
\nu_{M_0}(\{0\})=d(M_0^-)-d(M_0^+)=\delta_*(m_0^-)-\delta_*(m_0^+),
\]
where $\delta_*$ is any of the four functions $m\mapsto\dlow\{D_m>0\}$, $\dup\{D_m>0\}$, $\dlow\{D_m\ge0\}$, $\dup\{D_m\ge0\}$; all four are nonincreasing on $I-\frac12$ and have the same one-sided limits at interior points.

\textup{(c)} The set $\mathbb E=\{M\in I:\nu_M(\{0\})>0\}$ is countable, and $\sum_{M\in\mathbb E\cap[M_1,M_2]}\nu_M(\{0\})\le\bar d(M_1)-d(M_2)$ for $M_1<M_2$ in $I$.

\textup{(d)} For every $M=m+\frac12\in I\setminus\mathbb E$, the logarithmic densities of $\{D_m>0\}$, $\{D_m\ge0\}$ and $\{\sum_{p\le x}\chi^*(p)p^m\log p>0\}$ exist and equal $d(M)$, and for every function $\eps(x)\to0$ the near-ties $\{x:|D_m(x)|\le\eps(x)x^M/\log x\}$ have logarithmic density $0$.
\end{theorem}

\begin{proof}
(a) By \eqref{eq:transversal} and \eqref{eq:Gbounds}, $F_M/M\ge F_{M'}/M'+\kappa\lambda_\chi(M,M')$ a.e., so $\{F_{M'}\ge0\}\subset\{F_M>0\}$ up to a null set, and $d(M')\le\bar d(M')\le d(M)\le\bar d(M)$.

(b) $M\mapsto F_M\in L^2(\cA)$ is continuous (dominated convergence in the coefficients), so $M\mapsto\nu_M$ is weakly continuous, and by (I3) $d$ is lower and $\bar d$ upper semicontinuous. For $M\downarrow M_0$: $d(M)\le\bar d(M)\le d(M_0)$ by (a) and $\liminf d(M)\ge d(M_0)$, so $d(M_0^+)=\bar d(M_0^+)=d(M_0)$. For $M\uparrow M_0$: $\bar d(M)\ge d(M)\ge\bar d(M_0)$ and $\limsup\bar d(M)\le\bar d(M_0)$, so $\bar d(M_0^-)=d(M_0^-)=\bar d(M_0)$. Subtracting, $\nu_{M_0}(\{0\})=\bar d(M_0)-d(M_0)=d(M_0^-)-d(M_0^+)$. By \eqref{eq:sandwich} each $\delta_*$ lies between $d$ and $\bar d$, hence has the same one-sided limits, and is nonincreasing by (a).

(c) By (a), for $M_1\le M\le M_2$ the sets $S_M=\{F_M=0\}$ are pairwise disjoint mod null sets and contained in $\{F_{M_1}\ge0\}\setminus\{F_{M_2}>0\}$, of measure $\bar d(M_1)-d(M_2)$.

(d) If $\nu_M(\{0\})=0$ then $d(M)=\bar d(M)$ and \eqref{eq:sandwich} gives the densities; the $\theta$-form has the same limiting distribution (I1). Near-ties: $|D_m(x)|\le\eps(x)x^M/\log x$ iff $|E_M(x)|\le2M\eps(x)$; for fixed $\eps_0>0$ this set is eventually inside $\{|E_M|\le\eps_0\}$, of upper density $\le\nu_M([-\eps_0,\eps_0])$ by \eqref{eq:portmanteau}, which tends to $\nu_M(\{0\})=0$.
\end{proof}

\begin{corollary}[the classical race; GRH; \rd{2 reads}]\label{cor:classical}
Assume GRH for $L(s,\chi_{-4})$. Then $\nu_{1/2}(\{0\})=\delta(0^-)-\delta(0^+)$, where $\delta(m)$ is the lower (or upper) logarithmic density of $\{x:\sum_{p\le x}\chi^*(p)p^m>0\}$. If $m\mapsto\delta(m)$ is continuous at $m=0$, the logarithmic density of $\{x:\pi(x;4,3)>\pi(x;4,1)\}$ exists and near-ties have density $0$. For all but countably many $m\in(-\frac12,1.3239]$ the logarithmic density of $\{x:\sum_{p\le x}\chi^*(p)p^m>0\}$ exists. The same holds for $\chi_{-3}$ with $(-\frac12,1.6277]$.
\end{corollary}

The converse ``density exists $\Rightarrow$ no atom'' is not claimed: an atom at $0$ could in principle be carried by times at which $E_M$ is small and positive. The total mass of the atoms over $I=(0,M_\sigma]$ is $\le1-d(M_\sigma)<1$, since $\nu_M$ has mean $1$.

\begin{remark}[what Theorem~\ref{thm:jump} does not give]
The proof shows that $\mathbb E$ is countable, not that a given $m$ (in particular $m=0$) lies outside it. The family is real-analytic in $M$ ($M\mapsto F_M\in L^2(\cA)$ is), but analyticity of the family does not transfer to the distribution function without information on the law. The same countability holds along every one-parameter family $F+sH$ with $H$ continuous and positive on $\cA$, and for the same reason it never decides a prescribed parameter.
\end{remark}

\begin{theorem}[two-sided quantitative form; GRH, NCZ; \rd{2 reads}]\label{thm:twosided}
Let $M_1<M_0<M_2$ in $I$ and
\[
a_1=(M_0-M_1)M_0\lambda_\chi(M_1,M_0),\quad b_1=(M_0-M_1)\Bigl(\frac2{M_1}-M_0\lambda_\chi(M_1,M_0)\Bigr),
\]
\[
a_2=(M_2-M_0)M_0\lambda_\chi(M_0,M_2),\quad b_2=(M_2-M_0)\Bigl(\frac2{M_2}-M_0\lambda_\chi(M_0,M_2)\Bigr).
\]
Then $\nu_{M_0}((-a_1,a_2))\le d(M_1)-\bar d(M_2)\le\bar d(M_1)-d(M_2)\le\nu_{M_0}([-b_1,b_2])$. So $\nu_{M_0}((-t,t))=O(t)$ as $t\to0$ iff $d$ is Lipschitz at $M_0$, and $\nu_{M_0}((-t,t))\to0$ iff $d$ is continuous at $M_0$. Since a monotone function is differentiable almost everywhere, $\limsup_{t\to0}\nu_M((-t,t))/t<\infty$ for Lebesgue-almost every $M\in I$.
\end{theorem}

\begin{proof}
By \eqref{eq:transversal} and \eqref{eq:Gbounds}, a.e., $F_{M_1}/M_1-F_{M_0}/M_0\in[(M_0-M_1)\lambda_\chi(M_1,M_0),(M_0-M_1)(\frac2{M_1M_0}-\lambda_\chi(M_1,M_0))]$, and likewise for $(M_0,M_2)$. Hence a.e.: $F_{M_0}>-a_1\Rightarrow F_{M_1}>0$, and $F_{M_0}<a_2\Rightarrow F_{M_2}<0$; so $\{-a_1<F_{M_0}<a_2\}\subset\{F_{M_1}>0\}\setminus\{F_{M_2}\ge0\}$, a set of measure $d(M_1)-\bar d(M_2)$. Conversely $F_{M_1}\ge0\Rightarrow F_{M_0}\ge-b_1$ and $F_{M_2}\le0\Rightarrow F_{M_0}\le b_2$, so $\{F_{M_1}\ge0\}\setminus\{F_{M_2}>0\}\subset\{-b_1\le F_{M_0}\le b_2\}$. The rest follows with $M_{1,2}=M_0\mp t$ and continuity of $\lambda_\chi$.
\end{proof}

\begin{proposition}[the critical weight; GRH, NCZ; \rd{2 reads}]\label{prop:critical}
There is a version of the family such that for every $\theta$ outside a null set the map $M\mapsto\Phi_\theta(M):=F_M(\theta)/M$ is continuous and strictly decreasing on $I$, with $(M'-M)\lambda_\chi(M,M')\le\Phi_\theta(M)-\Phi_\theta(M')\le(M'-M)(\frac2{MM'}-\lambda_\chi(M,M'))$. Let $M^*(\theta)$ be its zero in $I$ (with the obvious conventions if it has none). Then for interior $M$: $d(M)=\hA(M^*>M)$, $\bar d(M)=\hA(M^*\ge M)$ and $\nu_M(\{0\})=\hA(M^*=M)$. So the exceptional weights are exactly the atoms of the law of the critical weight $M^*$.
\end{proposition}

\begin{proof}
Fix $M_1\in I$ and a Borel version of $F_{M_1}$; put $\Phi_\theta(M):=F_{M_1}(\theta)/M_1+\frac{M_1-M}{MM_1}G_{M,M_1}(\theta)$, with $G_{M,M'}:=G_{M',M}$ for $M>M'$ and $G_{M,M}:=1+\re\sum m_\gamma4M^2\theta_\gamma/(M+i\gamma)^2$. This is continuous in $M$ for every $\theta$, equals $F_M/M$ a.e.\ for each $M$ by \eqref{eq:transversal}, and the cocycle identity $\frac{M_1-M}{MM_1}G_{M,M_1}+\frac{M_2-M_1}{M_1M_2}G_{M_1,M_2}=\frac{M_2-M}{MM_2}G_{M,M_2}$ (an identity between absolutely convergent series) gives the two-sided bound for every $\theta$ from \eqref{eq:Gbounds}.
\end{proof}

\section{Time reversal and certified bounds}\label{sec:reversal}

Write $\theta_\gamma=e^{i\varphi_\gamma}$ and
\[
R:=\sum_{\gamma\in\Gamma}m_\gamma w_\gamma\cos\varphi_\gamma,\qquad Y:=\sum_{\gamma\in\Gamma}m_\gamma s_\gamma\sin\varphi_\gamma\in L^2(\cA).
\]
$R$ converges absolutely and uniformly (as $\sum m_\gamma/\gamma^2<\infty$), so it is continuous on $\cA$ and $|R|\le\sum m_\gamma w_\gamma=\sigma_M^2$ everywhere; note that $m_\gamma$ enters $R$ to the first power. Let $\mathrm{inv}(\theta)=\theta^{-1}$; it preserves $\hA$, since $\cA$ is a compact abelian group.

\begin{theorem}[time reversal; GRH, NCZ, no LI; \rd{2 reads}]\label{thm:reversal}
For every $M>0$:

\textup{(a)} $F_M=1+R/2+Y$ a.e., $R\circ\mathrm{inv}=R$ and $Y\circ\mathrm{inv}=-Y$. Hence
\[
F_M(\theta)+F_M(\theta^{-1})=2+R(\theta)\in[2-\sigma_M^2,\,2+\sigma_M^2]\qquad\text{for a.e. }\theta,
\]
and the law of $Y$ is symmetric about $0$.

\textup{(b)} For every real $t$: $\nu_M((-\infty,t])\le\nu_M([2-\sigma_M^2-t,\infty))$ and $\nu_M([t,\infty))\le\nu_M((-\infty,2+\sigma_M^2-t])$. Equivalently, there is a coupling $(X,X')$ of $\nu_M$ with itself such that $|X+X'-2|\le\sigma_M^2$ almost surely: $\nu_M$ is symmetric about its mean $1$ up to $\sigma_M^2$ in the $L^\infty$-Wasserstein distance.

\textup{(c)} If $\sigma_M^2<2$ (i.e.\ $M<M_\sigma(\chi)$), the sets $\{F_M\le0\}$ and $\mathrm{inv}\{F_M\le0\}$ are disjoint mod null sets, $2\nu_M((-\infty,0])=\hA\{|Y|\ge1+R/2\}$, and $\nu_M((-\infty,0])\le\frac12$. In particular the lower logarithmic density of $\{D_m>0\}$ is $\ge\frac12$ for every $-\frac12<m<M_\sigma(\chi)-\frac12$.
\end{theorem}

\begin{proof}
(a) $\re(r_\gamma\theta_\gamma)=\re r_\gamma\cos\varphi_\gamma-\im r_\gamma\sin\varphi_\gamma$, and under $\mathrm{inv}$, $\varphi_\gamma\mapsto-\varphi_\gamma$; the identities hold in $L^2(\cA)$ by comparing coefficients, hence a.e. (b) $X:=F_M$ and $X':=F_M\circ\mathrm{inv}$ have the same law, and $X'\le t$ implies $X\ge2-\sigma_M^2-t$. (c) $F_M\le0\iff Y\le-(1+R/2)$ and $F_M\circ\mathrm{inv}\le0\iff Y\ge1+R/2$, with $1+R/2\ge1-\sigma_M^2/2>0$; the two sets have equal measure. The density statement follows from \eqref{eq:sandwich}.
\end{proof}

Under LI the law of $F_M$ is exactly symmetric about $1$, by the independent rotations $\varphi_\gamma\mapsto\varphi_\gamma+\pi$; that symmetry is not available without LI. Devin \cite[Theorem~2.3]{Devin} obtains exact symmetry under the hypothesis that every integer relation $\sum k_\gamma\gamma=0$ has $\sum k_\gamma$ even, which puts the simultaneous half-turn in the hull. The involution used here, $y\mapsto-y$ on the orbit, is always available.

\begin{theorem}[GRH, NCZ, no LI, no non-resonance; \rd{2 reads}]\label{thm:cheb52}
If $\sigma_M^2<2$, then for every $0<c<2-\sigma_M^2$,
\[
\nu_M((-\infty,0])\le\frac{V_M+(1-c)^2}{c^2+(2-\sigma_M^2-c)^2}.
\]
In particular $\nu_M((-\infty,0])\le V_M/(1+(1-\sigma_M^2)^2)$ when $\sigma_M^2<1$ (take $c=1$); for $1\le\sigma_M^2<2$ the same bound holds because it is then weaker than Cantelli's $V_M/(1+V_M)$.
\end{theorem}

\begin{proof}
Let $N=\{F_M\le0\}$, $N'=\{F_M\ge2-\sigma_M^2\}$. By Theorem~\ref{thm:reversal}(a), $\mathrm{inv}N\subset N'$ mod null sets, so $\hA(N')\ge\hA(N)$, and $N\cap N'=\emptyset$. As $(F_M-c)^2\ge c^2$ on $N$ and $\ge(2-\sigma_M^2-c)^2$ on $N'$, $V_M+(1-c)^2=\int(F_M-c)^2\ge(c^2+(2-\sigma_M^2-c)^2)\hA(N)$. For $1\le\sigma_M^2<2$: $(1-\sigma_M^2)^2<1\le\sigma_M^2\le V_M$, so $V_M/(1+(1-\sigma_M^2)^2)>V_M/(1+V_M)$, and Cantelli's inequality (mean $1$, variance $V_M$) gives $\nu_M((-\infty,0])\le V_M/(1+V_M)$.
\end{proof}

For small $V_M$ this is about half of Cantelli's bound, as for a symmetric law. The next theorem adds finite arithmetic information about the lowest zeros. Let $\gamma_1<\dots<\gamma_N$ be the $N$ smallest ordinates and $T\ge3\gamma_N$ a height up to which the list of zeros is complete and all of them are simple; put $T'=T-2\gamma_N$ and $\mathcal K=\Gamma\cap(\gamma_N,T']$. For integer vectors $n$ indexed by $\{1,\dots,N\}$ write $n\cdot\gamma=\sum n_i\gamma_i$, $|n|_1=\sum|n_i|$. The \emph{non-resonance conditions} of \cite[\S6]{GO} are:
\begin{itemize}
\item[(NR1)] $n\cdot\gamma\ne0$ for all $n\ne0$ with $|n|_1\le4$;
\item[(NR2)] for $|n|_1\le3$, $|n\cdot\gamma|$ is not an ordinate outside $\{\gamma_1,\dots,\gamma_N\}$;
\item[(NR3)] for $1\le|n|_1\le2$, $n\cdot\gamma\ne\pm\gamma'\pm\gamma''$ for all ordinates $\gamma',\gamma''\in(\gamma_N,T]$.
\end{itemize}
These are finite statements, verified for the certified zero lists in exact integer arithmetic \cite{GO}.

\begin{theorem}[GRH, NCZ, finite non-resonance; \rd{2 reads}]\label{thm:nr}
Assume (NR1)--(NR3). Write $w_n,s_n$ for $w_{\gamma_n},s_{\gamma_n}$ and put
\[
\rho=\tfrac12\Bigl(\sigma_M^2-\sum_{n\le N}w_n\Bigr),\quad x_{\max}=\frac{\sum_{n\le N}w_n}{2(1-\rho)},\quad e=\frac{x_{\max}^2}{1-x_{\max}},\quad a_0=\frac{1+e}{1-\rho},\quad k_n=\frac{w_n}{2(1-\rho)^2},
\]
\[
V^o_{\mathcal K}=\tfrac12\sum_{\gamma\in\mathcal K}s_\gamma^2,\quad V^e=\tfrac18\sum_{\gamma\in\Gamma}m_\gamma^2w_\gamma^2,\quad V^o_t=V_M-V^e-\tfrac12\sum_{n\le N}s_n^2-V^o_{\mathcal K}.
\]
If $\rho<1$ and $x_{\max}<1$, then $2\nu_M((-\infty,0])\le J_R$, where
\[
J_R=\tfrac12a_0^2\sum s_n^2+\tfrac14\Bigl(\sum s_n^2\Bigr)\Bigl(\sum k_n^2\Bigr)-\tfrac18\sum s_n^2k_n^2+\Bigl(a_0^2+\tfrac12\sum k_n^2\Bigr)V^o_{\mathcal K}+\Bigl(a_0+\sum k_n\Bigr)^2V^o_t.
\]
$J_R$ is increasing in $V_M$ and decreasing in $V^e$, so a certified upper bound for $V_M$ and the lower bound $V^e\ge\frac18\sum_{\gamma\le T}w_\gamma^2$ may be used.
\end{theorem}

\begin{proof}
By the closed form of $\sigma_M^2$ and simplicity of the head zeros, $\rho=\frac12\sum_{\gamma\notin\text{head}}m_\gamma w_\gamma$, so with $R_h=\sum_{n\le N}w_n\cos\varphi_n$ we get $1+R/2\ge(1-\rho)+R_h/2=(1-\rho)(1+x)$, $x=R_h/(2(1-\rho))$, $|x|\le x_{\max}<1$. Since $1/(1+x)=1-x+x^2/(1+x)\le1-x+e$,
\[
\kappa:=\frac{1+e-x}{1-\rho}=a_0-\sum_{n\le N}k_n\cos\varphi_n\ge\frac1{(1-\rho)(1+x)}\ge\frac1{1+R/2}>0 .
\]
So $1+R/2>0$ everywhere, which is all that Theorem~\ref{thm:reversal}(c) uses, and $2\nu_M((-\infty,0])=\hA\{|Y|\ge1+R/2\}\le\hA\{\kappa|Y|\ge1\}\le\int\kappa^2Y^2$. Split $Y=Y_h+Y_{\mathcal K}+Y_t$ (head; $\mathcal K$; ordinates $>T'$ with multiplicities). $\kappa^2$ and $Y_h$ are trigonometric polynomials of degrees $2$ and $1$ in $e_{\gamma_1},\dots,e_{\gamma_N}$, and the integral of a monomial in the characters is $1$ if its total frequency is $0$ and $0$ otherwise. (1) $\int\kappa^2Y_h^2$ has degree $4$ in the head, so by (NR1) it equals its value for independent uniform phases; with $E\sin^2=E\cos^2=\frac12$, $E\sin^2\cos^2=\frac18$ this gives the first three terms. (2) $\int\kappa^2Y_hY_{\mathcal K}=\int\kappa^2Y_hY_t=0$: a head frequency $n\cdot\gamma$ with $|n|_1\le3$ would equal $\pm\gamma'$ for an ordinate outside the head, excluded by (NR2). (3) $\int\kappa^2Y_{\mathcal K}^2=(a_0^2+\frac12\sum k_n^2)V^o_{\mathcal K}$ by (NR3) and (NR1). (4) $\int\kappa^2Y_{\mathcal K}Y_t=0$ by (NR3) and the buffer $T'=T-2\gamma_N$: a relation $n\cdot\gamma\pm\gamma'\pm\gamma''=0$ with $\gamma''>T$ would force $\gamma''\le2\gamma_N+T'=T$. (5) $\int\kappa^2Y_t^2\le\|\kappa\|_\infty^2\int Y_t^2\le(a_0+\sum k_n)^2V^o_t$, since $\int Y_t^2=\frac12\sum_{\gamma>T'}m_\gamma^2s_\gamma^2=V_M-V^e-\frac12\sum_{n\le N}s_n^2-V^o_{\mathcal K}$.
\end{proof}

\begin{theorem}[Chebyshev's race, certified; GRH only; \rd{2 reads}]\label{thm:cert}
Assume GRH for $L(s,\chi_{-4})$. Then the lower logarithmic density of $\{x:\pi(x;4,3)>\pi(x;4,1)\}$ is
\[
>0.901150240\quad\text{(Theorem~\ref{thm:cheb52}, $c=1.019$)},\qquad>0.910924861\quad\text{(Theorem~\ref{thm:nr}, $N=20$)},
\]
and $\nu_{1/2}(\{0\})\le\nu_{1/2}((-\infty,0])<0.0890752$; so the upper and lower logarithmic densities of this set differ by less than $0.0891$. Assume GRH for $L(s,\chi_{-3})$. Then the lower logarithmic density of $\{x:\pi(x;3,2)>\pi(x;3,1)\}$ is $>0.929813746$ (Theorem~\ref{thm:cheb52}) and $>0.934761398$ (Theorem~\ref{thm:nr}, $N=20$), and $\nu_{1/2}(\{0\})<0.0652387$.
\end{theorem}

\begin{proof}
NCZ holds for $q=3,4$. Apply Theorems~\ref{thm:cheb52} and \ref{thm:nr} with the inputs certified in \cite{GO}: the zeros of $L(s,\chi_{-4})$ (resp.\ $L(s,\chi_{-3})$) are certified complete and simple up to $T=640.1657$ (resp.\ $700.0247$), as balls of radius $10^{-20}$ (sign changes of Hardy's function together with a winding-number count); $\sigma_{1/2}^2$ from the closed form ($0.155567979923586\ldots$, resp.\ $0.113229969857472\ldots$); the certified variance bound $V_{1/2}\le V^+=0.1696308345$ (resp.\ $0.1254514113$), from an explicit bound for the multiplicities of the zeros above $T$ \cite[Proposition~4.4]{GO}; and (NR1)--(NR3) for $N=20$, checked in exact integer arithmetic on the ordinates multiplied by $10^{24}$ (smallest margin $1.47\cdot10^{-6}$ for $q=4$, $7.94\cdot10^{-7}$ for $q=3$, against a tolerance of $4\cdot10^{-20}$). The evaluation is in Arb \cite{Arb} at $200$ bits (\file{density/reflection_certified.py}, output \file{density/reflection_certified_results.txt}); the set $\mathcal K$ is decided by certified comparisons. By \eqref{eq:sandwich} the lower density of $\{D_0>0\}$ is $\ge1-\nu_{1/2}((-\infty,0])$.
\end{proof}

Both reads re-ran \file{reflection_certified.py} with byte-identical output, and the second recomputed every number independently in mpmath at $40$ digits ($\nu\le0.08907513807014$, resp.\ $0.06523860156028$). Under GRH and LI the densities are $0.9959\ldots$ and $0.9990\ldots$ \cite{RS}. If one assumes in addition that all zeros up to height $10^8/q$ are simple (``tier B'', a by-product of Platt's verification of GRH \cite{Platt} that we have not seen stated as a theorem, \cite[Remark~4.6]{GO}), the non-resonance bound of \cite[Theorem~6.2]{GO} gives $0.922718033$ and $0.944555790$, which the present method does not improve ($0.919189115$, $0.941625186$). The two methods are complementary: the reflection bound wins when the unknown tail variance is the bottleneck (tier A near $m=0$), the one-sided multiplier of \cite{GO} when the head is well certified or the variance is tiny. Table~\ref{tab:weights} gives certified lower bounds for the lower logarithmic density of $\{D_m>0\}$ at other weights (tier A, values rounded down; ``prev.'' is the best previous value from \cite{GO}).

\begin{table}[ht]\small
\caption{Certified lower bounds for $\dlow\{D_m>0\}$, tier A (zeros certified up to $640$, resp.\ $700$).}\label{tab:weights}
\begin{tabular}{lccccc}\toprule
$m$&$-0.4$&$-0.25$&$0$&$0.25$&$0.5$\\\midrule
$\chi_{-4}$, Thm~\ref{thm:cheb52}&$0.99657$&$0.97790$&$0.90115$&$0.74212$&$0.47052$\\
$\chi_{-4}$, Thm~\ref{thm:nr}&$0.99659$&$\mathbf{0.97853}$&$\mathbf{0.91092}$&$\mathbf{0.78057}$&$0.52070$\\
$\chi_{-4}$, prev.&$\mathbf{0.9976}$&$0.9784$&$0.8842$&$0.7378$&$\mathbf{0.5983}$\\
$\chi_{-3}$, Thm~\ref{thm:nr}&&$\mathbf{0.98416}$&$\mathbf{0.93476}$&&$\mathbf{0.68600}$\\
$\chi_{-3}$, prev.&&$0.9830$&$0.9112$&&$0.6695$\\\bottomrule
\end{tabular}
\end{table}


\section{Sufficient conditions weaker than linear independence}\label{sec:sufficient}

In this section GRH$(\chi)$ and NCZ$(\chi)$ are assumed throughout. Each hypothesis below is a condition on the set $\Gamma$ (with multiplicities); none is known to hold for any $L$-function, and each is implied by LI (with simple zeros where stated). For $X\subset\RR$, $\spanQ X$ denotes the $\QQ$-span. In every case the conclusion ``no atoms'' gives, by \eqref{eq:portmanteau} and \eqref{eq:sandwich} exactly as in Theorem~\ref{thm:jump}(d): for every $m>-\frac12$ the logarithmic densities of $\{D_m>0\}$, $\{D_m\ge0\}$ and $\{x:D_m(x)>c\,x^M/\log x\}$ (every $c$) exist, and near-ties have logarithmic density $0$. In particular, for $\chi_{-4}$, the logarithmic density of $\{x:\pi(x;4,3)>\pi(x;4,1)\}$ exists.

\subsection{The tail does not span: TS, FD and tame directions}\label{sec:TS}

\begin{hypothesis}[TS]
There is $T>0$ with $\spanQ{\Gamma_{>T}}\ne\spanQ\Gamma$, i.e.\ some ordinate is not a rational linear combination of ordinates exceeding $T$. Equivalently, there is a $\QQ$-linear form $\varphi\ne0$ on $\spanQ\Gamma$ with $\varphi(\gamma)=0$ for all but finitely many $\gamma\in\Gamma$.
\end{hypothesis}

\begin{theorem}[GRH, NCZ, TS; \rd{1 read}]\label{thm:TS}
For every $M>0$ the measure $\nu_M$ is absolutely continuous.
\end{theorem}

\begin{proof}
Let $\varphi$ be as in TS and $\Lambda=\{\gamma\in\Gamma:\varphi(\gamma)\ne0\}$ (finite, nonempty); after scaling, $\varphi(G)\subset\ZZ$. Put $j(s)=(e^{i\varphi(\gamma)s})_{\gamma\in\Gamma}$. If $\sum n_\gamma\gamma=0$ then $\prod e^{i\varphi(\gamma)sn_\gamma}=1$, so $j(s)\in\cA$ by (I2), and $j:\RR\to\cA$ is a $2\pi$-periodic one-parameter subgroup with $e_g(j(s))=e^{i\varphi(g)s}$. Write $F_M=F'+P$ with $P=\re\sum_{\gamma\in\Lambda}m_\gamma r_\gamma e_\gamma$. The Fourier coefficients of $F'$ lie in $\ker\varphi$, so for every $s$, $F'(\theta j(s))=F'(\theta)$ for a.e.\ $\theta$; by Fubini there is a null set $N_1$ such that for $\theta\notin N_1$ this holds for a.e.\ $s\in[0,2\pi]$. For every $\theta$, $p_\theta(s):=P(\theta j(s))=\sum_{k\ne0}C_k(\theta)e^{iks}$ with
\[
C_k(\theta)=\tfrac12\sum_{\gamma\in\Lambda,\ \varphi(\gamma)=k}m_\gamma r_\gamma\theta_\gamma+\tfrac12\sum_{\gamma\in\Lambda,\ \varphi(\gamma)=-k}m_\gamma\bar r_\gamma\bar\theta_\gamma .
\]
Choose $k_0\ne0$ with $\{\gamma\in\Lambda:\varphi(\gamma)=\pm k_0\}\ne\emptyset$. Then $C_{k_0}$ is a trigonometric polynomial in pairwise distinct characters with nonzero coefficients, so its zero set $N_2$ is null (Lemma~\ref{lem:trigpoly}); for $\theta\notin N_2$, $p_\theta$ is non-constant and real-analytic. For a Lebesgue-null Borel $B$, by invariance and Fubini,
\[
\nu_M(B)=\int_\cA\frac1{2\pi}\mathrm{Leb}\{s\in[0,2\pi]:F_M(\theta j(s))\in B\}\,d\hA(\theta),
\]
and for $\theta\notin N_1\cup N_2$ the inner set equals $\{s:F'(\theta)+p_\theta(s)\in B\}$ up to a null set, which is null by Lemma~\ref{lem:analytic}.
\end{proof}

\begin{corollary}[what a non-absolutely-continuous law forces; GRH, NCZ; \rd{1 read}]\label{cor:TSatom}
If for some $M>0$ the measure $\nu_M$ is not absolutely continuous (in particular if it has an atom), then for every $T$ every ordinate of $L(s,\chi)$ is a rational linear combination of ordinates exceeding $T$.
\end{corollary}

\begin{example}[\rd{1 read}]\label{ex:TSweaker}
TS is strictly weaker than the hypothesis of Devin \cite[Theorem~2.2]{Devin} (a finite nonempty $\Lambda'$ with $\spanQ{\Lambda'}\cap\spanQ{\Gamma\setminus\Lambda'}=\{0\}$), as a condition on a set of frequencies. Let $b_0,b_1,\ldots>0$ be $\QQ$-independent with $b_i\to\infty$, and $X=\{b_0,b_0+b_1\}\cup\{b_i+b_j:1\le i<j\}$. The dual form $b_0^*$ is nonzero exactly on $b_0,b_0+b_1$, so TS holds. But for a finite $\Lambda'$ every $b_i$ ($i\ge1$) lies in $\spanQ{X\setminus\Lambda'}$ ($2b_i=(b_i+b_k)+(b_i+b_l)-(b_k+b_l)$ for large $k<l$), so the intersection is never $\{0\}$.
\end{example}

\begin{hypothesis}[FD, finite-level directions]
There are $N_0\ge1$ and $\eps\in(0,\frac12)$ such that for arbitrarily large $T$ there exist $\Lambda_T\subset\Gamma_T$ with $\#\Lambda_T\le N_0$ and $\gamma_T\in\Lambda_T$ with $\gamma_T\le T^{1/2-\eps}$ and $\gamma_T\notin\spanQ{\Gamma_T\setminus\Lambda_T}$.
\end{hypothesis}

TS implies FD (take the fixed set $\Lambda=\{\gamma:\varphi(\gamma)\ne0\}$).

\begin{theorem}[GRH, NCZ, FD; \rd{1 read}]\label{thm:FD}
For every $M>0$ the measure $\nu_M$ has no atoms. If FD holds for all $T\ge T_0$, then $\sup_c\nu_M([c-\ell,c+\ell])\ll\ell^\beta$ as $\ell\to0$, with $\beta=\eps/(4N_0(1-\eps))$.
\end{theorem}

\begin{proof}
Fix $M$ and $T$ as in FD; there is a homomorphism $\varphi:G_T\to\ZZ$ ($G_T$ generated by $\Gamma_T$) vanishing on $\Gamma_T\setminus\Lambda$ and nonzero at $\gamma_T$, where $\Lambda=\Lambda_T$, $n=\#\Lambda\le N_0$. Let $\cA_T$ be the closure of $\{(e^{i\gamma y})_{\gamma\in\Gamma_T}\}$ and $\pi_T:\cA\to\cA_T$ the restriction (onto, measure preserving). By (I7),
\[
\|F_M-F_{M,T}\circ\pi_T\|_2^2=V_M^{>T}\le C_1M^2\log^2(T+3)/T.
\]
As in Theorem~\ref{thm:TS}, $j(s)=(e^{i\varphi(\gamma)s})_{\gamma\in\Gamma_T}$ is a $2\pi$-periodic subgroup of $\cA_T$ and $F_{M,T}(\theta j(s))=F'(\theta)+\sum_{k\in K}C_k(\theta)e^{iks}$, $\#K\le2n$, $0\notin K$. Put $k_*=\varphi(\gamma_T)$ and $a=\frac12m_{\gamma_T}|r_{\gamma_T}|\ge2M/(M+T^{1/2-\eps})$. Fix $c$, $\ell>0$, $\theta$, let $q(s)=F'(\theta)-c+\sum_KC_k(\theta)e^{iks}$ and $E_\theta=\{s\in[0,2\pi]:|q(s)|\le\ell\}$. Since $\sup|q|\ge|C_{k_*}(\theta)|$, (I6) gives $|E_\theta|/2\pi\le A_N(\ell/|C_{k_*}(\theta)|)^{1/(2N_0)}$ when $\ell\le|C_{k_*}(\theta)|$. Next, $\hA\{|C_{k_*}|\le\eta a\}\le A_N\eta^{1/N_0}$ for $0<\eta<1$: if $C_{k_*}$ has one term the set is empty; otherwise choose $\psi:G_T\to\ZZ$ injective on the finitely many frequencies of $C_{k_*}$, move along $j_\psi$, and apply (I6) to a trigonometric polynomial with at most $N_0$ distinct integer frequencies, one coefficient of modulus $a$. By invariance and Fubini,
\[
\nu_M([c-\tfrac\ell2,c+\tfrac\ell2])\le A_N\eta^{1/N_0}+A_N\Bigl(\frac\ell{\eta a}\Bigr)^{1/(2N_0)}+\frac{4C_1M^2\log^2(T+3)}{T\ell^2}\qquad(\ell\le\eta a).
\]
Take $\ell=T^{-1/2+\eps/2}$ and $\eta=T^{-\eps/4}$; then $\ell/(\eta a)\le T^{-\eps/4}(M+1)/(2M)$, and the right side is $\ll T^{-\eps/(8N_0)}+T^{-\eps}\log^2(T+3)$, uniformly in $c$. Let $T\to\infty$ along FD. If every $T\ge T_0$ is allowed, take $T$ with $T^{-1/2+\eps/2}=2\ell_1$.
\end{proof}

\emph{Comparison with the literature} (statements read in the arXiv versions; \file{density/LITERATURE.md}). One self-sufficient ordinate ($\gamma\notin\spanQ{\Gamma\setminus\{\gamma\}}$; Devin \cite[Theorem~5.4(2)]{Devin}, Martin--Ng \cite[Definition~1.3]{MN}) implies Devin's direct-sum hypothesis \cite[Theorem~2.2]{Devin} (also Bailleul \cite[Theorem~1.7]{Bailleul}, Devin \cite[Corollary~2]{Devin2}), which implies TS, which implies FD. Devin's Theorem~5.4(1) (for every $T\ge T_0$ some $\gamma_T\le T^{1/2-\eps}$ not in $\spanQ{\Gamma_T\setminus\{\gamma_T\}}$) is FD with $N_0=1$ and ``every $T\ge T_0$''. TS does not imply it (Example~\ref{ex:TSweaker}) and it does not imply TS (an example supplied by the first referee: \file{density/DENSITY.md}, Example~3.3a; it was rewritten after the second read, following that read's suggestion, and not re-read). So FD is implied by each of these hypotheses and strictly weaker than each. Theorem~\ref{thm:TS} gives absolute continuity under TS, whereas \cite[Theorem~2.2]{Devin} gives continuity under the direct-sum hypothesis. The gain is modest.

\begin{definition}
A homomorphism $\varphi:G\to\RR$ is \emph{tame} if $|\varphi(\gamma)|\le C_\varphi\log(\gamma+3)$ for all $\gamma\in\Gamma$. For $v\in\RR$ let $C^\varphi_v=\sum_{g\in\pm\Gamma,\ \varphi(g)=v}m_g\frac{2M}{M+ig}e_g\in L^2(\cA)$ (with $m_{-\gamma}=m_\gamma$), and $j_\varphi(s)=(e^{i\varphi(\gamma)s})_\gamma$, a one-parameter subgroup of $\cA$ by (I2).
\end{definition}

\begin{theorem}[tame directions; GRH, NCZ; \rd{1 read}]\label{thm:tame}
Let $\varphi$ be tame. \textup{(a)} For a.e.\ $\theta$ there is a real-analytic $h_\theta$ with $F_M(\theta j_\varphi(s))=h_\theta(s)$ for a.e.\ $s$. \textup{(b)} If $S=\{F_M=c\}$ has positive measure, then $S$ is invariant under $j_\varphi(\RR)$ up to null sets, $C^\varphi_v=0$ a.e.\ on $S$ for every $v\ne0$, and $1+C^\varphi_0=c$ a.e.\ on $S$. \textup{(c)} If for some $v\ne0$ the set $\{g\in\pm\Gamma:\varphi(g)=v\}$ is finite and nonempty, then $\nu_M$ is absolutely continuous.
\end{theorem}

\begin{proof}[Proof (outline; complete proof in \file{density/DENSITY.md}, Theorem~3.5)]
(a) With $c_0=1/(4C_\varphi)$, the partial sums $F^z_N=1+\sum_{|g|\le N}m_g\frac{2M}{M+ig}e^{i\varphi(g)z}\theta_g$ are entire in $z$, and $|e^{i\varphi(g)z}|\le(|g|+3)^{1/4}$ on $|\im z|\le c_0$; by orthonormality and (I7) they are $L^2$-Cauchy uniformly on the strip, and the sub-mean-value property gives a.e.\ locally uniform convergence along a subsequence to a holomorphic $H_\theta$, with $F^s_N(\theta)=F^0_N(\theta j_\varphi(s))$. (b) For a.e.\ $\theta\in S$, $h_\theta=c$ on a set of positive measure, hence everywhere; so $S$ is $j_\varphi$-invariant, and the projections $P_v=\lim\frac1R\int_0^Re^{-ivs}U_s\,ds$ (mean ergodic theorem along $j_\varphi$), which commute with multiplication by $\mathbf 1_S$, applied to $\mathbf 1_SF_M=c\mathbf 1_S$ give the claims. (c) If $\nu_M(B)>0$ for a null Borel $B$, the set where $h_\theta$ is constant has positive measure, and the argument of (b) gives $C^\varphi_v=0$ on it, contradicting Lemma~\ref{lem:trigpoly} for the nonzero trigonometric polynomial $C^\varphi_v$.
\end{proof}

\begin{corollary}[necessary conditions for an atom; GRH, NCZ; \rd{1 read}]\label{cor:tame}
If $\nu_M(\{c\})>0$ and $S=\{F_M=c\}$, then for every tame $\varphi$: $\mathbf 1_S$ is $j_\varphi(\RR)$-invariant; every fibre $\{g\in\pm\Gamma:\varphi(g)=v\}$, $v\ne0$, is empty or infinite; and if $c\ne1$ the fibre over $0$ is infinite.
\end{corollary}

\subsection{The tail does not generate: TS\texorpdfstring{$_\ZZ$}{\_Z}}\label{sec:TSZ}

\begin{proposition}[finite-index criterion; GRH, NCZ; \rd{2 reads}]\label{prop:index}
Let $H\subset G$ be a subgroup of finite index $N$ such that $\Lambda=\Gamma\setminus H$ is finite. Then $\nu_M(\{c\})\le1/N$ for every $M>0$ and $c\in\RR$.
\end{proposition}

\begin{proof}
Let $K=H^\perp=\{\theta\in\cA:e_h(\theta)=1\ \forall h\in H\}$, a subgroup of order $N$ (dual of $G/H$ by (I2)). Since $G=H+\langle\Lambda\rangle$ and the characters separate points, for $k\ne k'$ in $K$ some $\gamma\in\Lambda$ has $e_\gamma(k)\ne e_\gamma(k')$. Write $F_M=F_H+F_\Lambda$ with $F_\Lambda=\re\sum_{\gamma\in\Lambda}m_\gamma r_\gamma e_\gamma$; the coefficients of $F_H$ lie in $H$, so $F_H(\cdot k)=F_H$ a.e., and after replacing $F_H$ by its average over $K$ this holds everywhere. Let $E=\{F_M=c\}$ and $S(\theta)=\{k\in K:\theta k\in E\}$; then $\int|S(\theta)|\,d\hA=N\nu_M(\{c\})$. If $k\ne k'$ both lie in $S(\theta)$, then $F_\Lambda(\theta k)=F_\Lambda(\theta k')$, so $\theta$ is a zero of the trigonometric polynomial $\re\sum_{\gamma\in\Lambda}m_\gamma r_\gamma(e_\gamma(k)-e_\gamma(k'))e_\gamma$, which has a nonzero coefficient; by Lemma~\ref{lem:trigpoly} this happens on a null set. So $|S(\theta)|\le1$ a.e.
\end{proof}

\begin{hypothesis}[TS$_\ZZ$]
$[G:\langle\Gamma_{>T}\rangle]\to\infty$ as $T\to\infty$ (the index is nondecreasing, and counts as $\infty$ when infinite).
\end{hypothesis}

\begin{theorem}[GRH, NCZ, TS$_\ZZ$; \rd{2 reads}]\label{thm:TSZ}
For every $M>0$ the measure $\nu_M$ has no atoms.
\end{theorem}

\begin{proof}
If some index $[G:\langle\Gamma_{>T}\rangle]$ is infinite, TS holds: otherwise every $\lambda\in\Gamma_T$ has a positive multiple in $\langle\Gamma_{>T}\rangle$ (as $\spanQ{\Gamma_{>T}}=\spanQ\Gamma$), and $G/\langle\Gamma_{>T}\rangle$ is generated by finitely many elements of finite order. Then Theorem~\ref{thm:TS} applies. Otherwise Proposition~\ref{prop:index} with $H=\langle\Gamma_{>T}\rangle$ gives $\nu_M(\{c\})\le1/[G:H]\to0$.
\end{proof}

\begin{corollary}[what an atom forces; GRH, NCZ; \rd{2 reads}]\label{cor:TSZatom}
If $\nu_M(\{c\})=\alpha>0$, then $\langle\Gamma_{>T}\rangle$ is eventually constant, equal to some $H_\infty$ with $[G:H_\infty]\le1/\alpha$, and for every $T$ each ordinate $\gamma$ has a multiple $n\gamma$, $1\le n\le1/\alpha$, that is an integer linear combination of ordinates exceeding $T$.
\end{corollary}

TS implies TS$_\ZZ$, and TS$_\ZZ$ is strictly weaker: $X=\{2^k\alpha:k\ge0\}$ has $\langle X_{>T}\rangle=2^{k(T)}\ZZ\alpha$, of index $\to\infty$, while TS fails. TS$_\ZZ$ and FD are incomparable ($X$ fails FD; the set of \file{density/DENSITY.md}, Example~3.3a, satisfies FD and fails TS$_\ZZ$). Devin's parity hypothesis \cite[Theorem~2.3]{Devin} implies neither ($X=\{(2k+1)\alpha:k\ge0\}$). An atom of mass $\alpha$ makes $G/H_\infty$ a finite group of order $\le1/\alpha$; it does not force it to be trivial, and Corollary~\ref{cor:TSZatom} gives no contradiction with the certified bound of Theorem~\ref{thm:cert}.

\subsection{A low vertex: hypothesis V}\label{sec:V}
For a trigonometric polynomial $P=\sum_{g\in S}c_ge_g$ on the dual of a torsion-free group ($c_g\ne0$), a point $g_0\in S$ is a \emph{vertex} if $g_0\notin\mathrm{conv}(S\setminus\{g_0\})$ in $G\otimes\RR$; equivalently some homomorphism $\varphi:\langle S\rangle\to\ZZ$ has $\varphi(g_0)>\varphi(g)$ for $g\in S\setminus\{g_0\}$. Write $\mathfrak m(f)=\int\log|f|$ (the logarithmic Mahler measure).

\begin{lemma}[Jensen--Newton; classical; \rd{1 read}]\label{lem:JN}
If $g_0$ is a vertex of $\mathrm{conv}(S)$, then $\mathfrak m(P)\ge\log|c_{g_0}|$.
\end{lemma}

\begin{proof}
Reduce to the torus dual to $\langle S\rangle$ and choose $\varphi$ as above. For each $\theta$, $s\mapsto P(\theta j(s))$ equals $e^{ik_-s}p_\theta(e^{is})$ for a polynomial $p_\theta$ of degree $k_+-k_-$ whose leading coefficient is $c_{g_0}e_{g_0}(\theta)$ ($k_\pm$ the extreme values of $\varphi$ on $S$). By Jensen's formula, $\frac1{2\pi}\int_0^{2\pi}\log|p_\theta(e^{is})|ds\ge\log|c_{g_0}|$ for every $\theta$; integrate over $\theta$ (Tonelli, separately for $\log^\pm$). (The lemma and this proof are classical; for $\gamma_0$ maximal in a group order, part (a) of Theorem~\ref{thm:V} below also follows from Jensen's inequality for $H^1$ of a compact group with ordered dual, Helson--Lowdenslager \cite{HL}.)
\end{proof}

\begin{lemma}[small ball from the Mahler measure; \rd{1 read}]\label{lem:smallball}
On a probability space, for $f\in L^2$ with $\mathfrak m(f)>-\infty$ and $0<\delta<1$: $\mu\{|f|<\delta\}\log(1/\delta)\le-\mathfrak m(f)+\frac12\log^+\int|f|^2+\frac1{2e}$.
\end{lemma}

\begin{proof}
With $B=\{|f|<\delta\}$, $\eta=\mu(B)$: $\int_B\log|f|\le\eta\log\delta$, and Jensen's inequality on $B^c$ gives $\int_{B^c}\log|f|\le\frac12\log^+\int|f|^2+\frac{1-\eta}2\log\frac1{1-\eta}$; and $x\log(1/x)\le1/e$.
\end{proof}

\begin{hypothesis}[V$(\gamma_0)$]
For $\gamma_0\in\Gamma$: for every $T\ge\gamma_0$, $\gamma_0$ is a vertex of $\mathrm{conv}(\pm\Gamma_T)$. Equivalently \textup{(\file{density3/MAHLER.md}, Proposition~3.1)}: every nontrivial integer relation $\sum k_\gamma\gamma=0$ among the distinct ordinates satisfies $|k_{\gamma_0}|<\sum_{\gamma\ne\gamma_0}|k_\gamma|$; equivalently, some total group order on $G$ has $\gamma_0=\max(\pm\Gamma)$.
\end{hypothesis}

\begin{theorem}[GRH, NCZ, V$(\gamma_0)$; \rd{1 read}]\label{thm:V}
For every $M>0$: \textup{(a)} $\int\log|x-c|\,d\nu_M(x)\ge\log(2Mm_{\gamma_0}/|M+i\gamma_0|)$ for every $c\in\CC$; \textup{(b)} for $c\in\RR$, $0<\eps<1$,
\[
\nu_M((c-\eps,c+\eps))\le\frac{\log\frac{|M+i\gamma_0|}{2Mm_{\gamma_0}}+\frac12\log^+((1-c)^2+V_M)+\frac1{2e}}{\log(1/\eps)};
\]
\textup{(c)} $\nu_M$ is absolutely continuous.
\end{theorem}

\begin{proof}
(a) For $T\ge\gamma_0$, the support of $F_{M,T}-c$ is $\pm\Gamma_T$ together with $0$, the midpoint of $\pm\gamma_0$, so $\gamma_0$ is a vertex and Lemma~\ref{lem:JN} gives $\mathfrak m(F_{M,T}-c)\ge\log(2Mm_{\gamma_0}/|M+i\gamma_0|)$. As $F_{M,T}\circ\pi_T\to F_M$ in $L^2$, reverse Fatou applied to $|f_T|-\log|f_T|\ge0$ along an a.e.-convergent subsequence gives the bound for $F_M$. (b) Lemma~\ref{lem:smallball}. (c) A compactness argument (inverse limit of the nonempty compact sets $\{\varphi\in\mathrm{Hom}(G_T,\RR):\varphi(\gamma_0)=1,\ |\varphi(\gamma)|\le1\ \forall\gamma\in\Gamma_T\}$) gives a homomorphism $\varphi:G\to\RR$ with $\varphi(\gamma_0)=1$ and $|\varphi|\le1$ on $\Gamma$; it is tame. As in (a), $\mathfrak m(C^\varphi_1)>-\infty$, so $C^\varphi_1\ne0$ a.e.; the proof of Theorem~\ref{thm:tame}(c) then applies, with Jensen in place of Lemma~\ref{lem:trigpoly}.
\end{proof}

If V$(\gamma_1)$ holds for $\gamma_1=\gamma_1(\chi_{-4})=6.0209\ldots$, then at $M=\frac12$, with simple zeros, the constant in (b) at $c=0$ is $2.0549$ (numerical); e.g.\ $\nu_{1/2}((-\eps,\eps))\le0.0892$ at $\eps=10^{-10}$, which is weaker than the unconditional bound of Theorem~\ref{thm:cert} (the point of (b) is the decay as $\eps\to0$). TS implies V$(\gamma_0)$ for some $\gamma_0$, and V is strictly weaker: $\Gamma_0=\{b+nd:n\ge0\}$ with $b/d\notin\QQ$ satisfies V$(b)$ but not TS, not FD, and has no tame direction with a finite nonzero fibre (\file{density3/MAHLER.md}, Proposition~3.6, Example~3.7). FD does not imply V (Example~3.8 there, given with less detail). V for any single zero implies that two ordinates of $L(s,\chi)$ have an irrational ratio, the $L(s,\chi)$-analogue of item~1 of \cite{HKMN}, which is open; like TS and FD it cannot be verified by a finite computation.

\begin{theorem}[low vertices; GRH, NCZ; \rd{1 read}]\label{thm:LV}
Let $\mu_T$ be the smallest $\gamma\in\Gamma_T$ that is a vertex of $\mathrm{conv}(\pm\Gamma_T)$ (the largest element of $\Gamma_T$ always is). Then for every $M>0$ and $c\in\RR$,
\[
\nu_M(\{c\})\le2\liminf_{T\to\infty}\frac{\log^+\mu_T}{\log T}.
\]
\end{theorem}

\begin{proof}[Proof (outline; \file{density3/MAHLER.md}, \S4)]
With $\delta_T(c)=-\mathfrak m(F_{M,T}-c)$ and $\tau_T=\log(1/V_M^{>T})=\log T+O(\log\log T)$: if $\nu_M(\{c\})=\eta>0$, then on $S=\{F_M=c\}$ one has $F_{M,T}\circ\pi_T-c=-(F_M-F_{M,T}\circ\pi_T)$, and concavity of $\log$ on $S$ and $S^c$ gives $\delta_T(c)\ge\frac\eta2\tau_T-O(1)$ for every $T$. Lemma~\ref{lem:JN} at the vertex $\mu_T$ gives $\delta_T(c)\le\log(|M+i\mu_T|/2M)$.
\end{proof}

So ``low vertices'' ($\log\mu_T=o(\log T)$ along a sequence) exclude atoms. In rank one with a dense spectrum the only vertex is the top ordinate, and the bound is void; all the atom models of Section~\ref{sec:barriers} are of this kind.

\subsection{Graded linear independence}\label{sec:GLI}
For a finitely supported $k\in\ZZ^{(\Gamma)}$ the \emph{length} is $\|k\|_1$, and $k$ is a \emph{relation} if $\sum k_\gamma\gamma=0$. For $L\ge1$ let $S_L$ be the set of ordinates occurring in some nonzero relation of length $\le L$.

\begin{lemma}[moment comparison; no hypothesis; \rd{2 reads}]\label{lem:moment}
Let $N\ge1$ and $W,Z$ real random variables with $EZ^j=EW^j$ for $0\le j\le2N$. Suppose $W$ has a density, and $W=\mu+U+V$ with $|U|\le b$ a.s.\ and $E|V|^{2N}\le2\cdot N!\cdot(2\sigma^2)^N$. For $x\in\RR$ put $R=16(|x-\mu|+b+2\sqrt N\sigma)$ and $D_x=\operatorname{ess\,sup}_{|t-x|\le R}f_W(t)$. Then $P(|Z-x|\le2R/(N+1))\le352D_xR/(N+1)+56\cdot16^{-N}$.
\end{lemma}

\begin{proof}[Proof (outline; \file{density2/INVLO.md}, Lemma~A)]
Test against $p(t)=q((t-x)/R)^2$, where $q$ is a ``needle'' polynomial of degree $\le N$ built from the Fej\'er kernel: $q(0)=1$, $0\le q\le2$ on $[-1,1]$, $\int_{-1}^1q^2\le8\pi/(N+1)$, $q\ge0.27$ on $|s|\le2/(N+1)$, and $|q(s)|\le2(2|s|)^N$ for $|s|\ge1$. Then $P(|Z-x|\le2R/(N+1))\le14Ep(Z)=14Ep(W)$, and $Ep(W)$ is bounded using the density near $x$ and the moment bound far away.
\end{proof}

The rate $N^{-1/2}$ is sharp (Gauss--Hermite quadrature of a Gaussian). The following is a form of the inverse Littlewood--Offord principle (cf.\ \cite{NV}) on a subtorus: \emph{near-atoms force short integer relations}.

\begin{proposition}[finite models; no hypothesis; \rd{2 reads}]\label{prop:Bprime}
Let $\Gamma'\subset(0,\infty)$ be finite, $c_\gamma\in\CC$, $F=1+\re\sum_{\gamma\in\Gamma'}c_\gamma e_\gamma$ on the closure torus of $\{(e^{i\gamma y})_{\gamma\in\Gamma'}\}$, $L\ge2$, $N=\lfloor L/2\rfloor$, $B=\sum_{\gamma\in S_L}|c_\gamma|$, $\sigma_Y^2=\sum_{\gamma\in\Gamma'\setminus S_L}|c_\gamma|^2$, and let $D$ bound the density of $Y=\re\sum_{\gamma\in\Gamma'\setminus S_L}c_\gamma e^{i\theta_\gamma}$ with independent uniform phases. Then, with $R=16(|x-1|+B+2\sqrt N\sigma_Y)$, $\lambda\{|F-x|\le2R/(N+1)\}\le352DR/(N+1)+56\cdot16^{-N}$.
\end{proposition}

\begin{proof}
Let $W=1+X+Y$ with $X=\re\sum_{\gamma\in S_L}c_\gamma e_\gamma$ (Haar on its own torus) independent of $Y$. Expanding $(F-1)^j$ and $(W-1)^j$ into sequences of characters, a sequence contributes iff its net vector $k$ is a relation (resp.\ a relation supported on $S_L$); for $j\le L$ these agree, as nonzero relations of length $\le L$ are supported in $S_L$. So the first $L\ge2N$ moments agree, and Lemma~\ref{lem:moment} applies with $\mu=1$, $U=X$, $V=Y$ (Hoeffding).
\end{proof}

\begin{hypothesis}[GLI$(\psi)$]
$\psi:[0,\infty)\to[0,\infty)$ is nondecreasing, and every nonzero relation $k$ has $\|k\|_1\ge\psi(\max\{\gamma:k_\gamma\ne0\})$.
\end{hypothesis}

\begin{theorem}[GRH, NCZ, GLI$(\psi)$ with $\psi(t)\log(t+3)/t\to\infty$; \rd{2 reads}]\label{thm:GLI}
For every $M>0$ the measure $\nu_M$ has no atoms.
\end{theorem}

\begin{proof}[Proof (outline; \file{density2/INVLO.md}, Theorem~C)]
Let $T_L=\sup\{t:\psi(t)\le L\}$; then $S_L\subset\Gamma_{T_L}$, so $B\le C_\chi M\log^2(T_L+3)$ (here and below $C_\chi$, $C$ denote constants depending on $\chi$, resp.\ on $c,M,\chi$). The block $I=\Gamma\cap(T_L,2T_L]$ lies outside $S_L$ and contains $\gg T_L$ distinct ordinates (Riemann--von Mangoldt; zeros are not assumed simple), with coefficients between $M/(2T_L)$ and $4MC_\chi\log(2T_L+3)/T_L$; a density bound for random cosine sums (product of Bessel functions $J_0$) gives $D\le C(1+(T_L/\log T_L)^{1/2})$. Proposition~\ref{prop:Bprime} for the truncation $F_{M,T}$ (with $T\ge T_L$, $N=\lfloor L/2\rfloor$, $\rho=2R/(N+1)$), plus the truncation error $\le4V_M^{>T}/\rho^2$ and then $T\to\infty$, gives $\nu_M(\{c\})\le C(1+(T_L/\log T_L)^{1/2})[(1+\log^2T_L)/L+L^{-1/2}]+o(1)$, and $\psi(t)\ge Kt/\log t$ for large $t$ makes this $\ll K^{-1/2}$.
\end{proof}

\begin{corollary}[what an atom forces; GRH, NCZ; \rd{2 reads}]\label{cor:GLIatom}
If $\nu_M(\{c\})>0$, there is $C<\infty$ such that infinitely many nonzero relations $k$ satisfy $\|k\|_1\le C\,t_k/\log(t_k+3)$, where $t_k$ is the largest ordinate involved. More strongly (\file{density2/INVLO.md}, Proposition~E): if $\nu_M(\{c\})=\alpha$, then for every $\delta\in(0,1]$, every large $L$ and every $T$ with $\log^3L\le T\le C_2^{-1}\alpha^2\delta L\log L$, all but $\delta T\log T$ of the zeros in $(T,2T]$ (with multiplicity) lie in relations of length $\le L$.
\end{corollary}

LI implies GLI$(\psi)$ for every $\psi$, and GLI is incomparable with TS, FD and TS$_\ZZ$ (\file{density2/INVLO.md}, \S3; e.g.\ $\{\alpha u_n/2^{e_n}\}$ with distinct odd primes $u_n$ and rapidly increasing $e_n$ satisfies GLI and fails TS, FD, TS$_\ZZ$). Proposition~\ref{prop:Bprime} is effective, but its constants are far too large to improve Theorem~\ref{thm:cert} from a finite search for relations.

\subsection{Prime-side hypotheses}\label{sec:prime}
Fix $v\in C_c^\infty(-\frac12,\frac12)$ even, $v\ge0$, $\int v=1$, and $w=v*v$; then $\hat w=\hat v^2\ge0$ is even and Schwartz, $\hat w(0)=1$, $\hat w>0$ on $(-\pi,\pi)$; put $c_w=\min_{|\xi|\le1}\hat w(\xi)>0$. For $T>0$ the \emph{windowed prime sum} is $\Pi_T(y)=\sum_n\chi(n)\Lambda(n)n^{-1/2}\,Tw(T(y-\log n))$, a weighted sum over a window of relative width $\asymp1/T$ around $x=e^y$. By the explicit formula, $\Pi_T(y)=\Omega_T(y)-G_T(\iota(y))$ with $G_T=\sum_\gamma m_\gamma\hat w(\gamma/T)(e_\gamma+e_{-\gamma})$ and $\Omega_T(y)\to0$ (numerically confirmed for $\chi_{-4}$ to $5.4\cdot10^{-4}$). So the log-averaged moments $M_k(T)=\lim_Y\frac1Y\int^Y\Pi_T(y)^kdy$ exist and equal $(-1)^k\int G_T^k\,d\hA$. List the zeros with positive ordinate with multiplicity as $Z^+$, put $a_\rho=\hat w(\gamma_\rho/T)$, and call a nonzero $\kappa\in\ZZ^{(Z^+)}$ with $\sum\kappa_\rho\gamma_\rho=0$ a \emph{multiset relation} (two copies of a multiple zero give one of length $2$). Let $\mu_k(T)=E[(\sum_\rho2a_\rho\cos\theta_\rho)^k]$ with independent uniform phases.

\begin{theorem}[duality; GRH, NCZ; \rd{2 reads}]\label{thm:P1}
For all $T>0$, $k\ge1$: $(-1)^kM_k(T)=\mu_k(T)+\Delta_k(T)$, where $\Delta_k(T)=\sum'\prod_{i\le k}a_{\rho_i}\ge0$, the sum over sequences $(\rho_i,\sigma_i)\in(Z^+\times\{\pm1\})^k$ whose net vector is a multiset relation. In particular odd moments are $\le0$; for each $L$, $\Delta_\ell(T)=0$ for all $\ell\le L$ iff there is no multiset relation of length $\le L$ among the $\rho$ with $a_\rho>0$; each relation of length $\ell$ among zeros of height $\le T$ contributes $\ge c_w^\ell$ to $\Delta_\ell(T)$; and $\mu_k(T)$ is a polynomial in finite prime sums (explicit formula at $y=0$).
\end{theorem}

\begin{proof}
Expand $E_{\hA}[G_T^k]$; $E[e_g]=\mathbf 1_{g=0}$, while in the independent model a sequence contributes iff its net vector is $0$. The difference is the sum over sequences with nonzero net vector and total frequency $0$, all terms $\ge0$. The other statements follow by realising a relation of length $\ell$ by a sequence of length $\ell$, padding with cancelling pairs, and the explicit formula for $\sum_\rho\hat w(\gamma_\rho/T)^{2j}$.
\end{proof}

So at this level a prime-side correlation hypothesis is a zero-side relation hypothesis in disguise. Its quantitative form (moments up to order $L$ match the prediction to within $c_w^\ell$ at one scale $T\ge C_1L^{3/2}\log^2L$) bounds every atom by $C(1+|c-1|)L^{-1/2}$, via Proposition~\ref{prop:Bprime} (\file{density2/PRIMECORR.md}, Theorem~P2, \rd{2 reads}); the infinite version is equivalent to LI with simple zeros. Mixed moments of two windows with disjoint frequency bands do better.

\emph{Band decoupling.} Let $L\ge2$, $T^*\ge1$, a band $0<s_1<s_2<T^*$ and a margin $g\in(0,1)$. Take even windows $\beta_I,\beta_R\in C_c^\infty(\RR)$ with values in $[0,1]$, $\beta_I=1$ on $[s_1+g,s_2-g]$ and $\operatorname{supp}\beta_I\cap[0,\infty)\subset[s_1,s_2]$, $\beta_R=1$ on $[0,s_1-g]\cup[s_2+g,T^*]$ and $\operatorname{supp}\beta_R\cap[0,\infty)\subset[0,s_1]\cup[s_2,T^*+1]$. With $b=\check\beta$, the limits $\Pi_\beta(y)=\lim_X\sum_{n\le X}\chi(n)\Lambda(n)n^{-1/2}b(y-\log n)$ exist and $\Pi_\beta(y)=-G_\beta(\iota(y))+o(1)$, $G_\beta=\sum_\gamma m_\gamma\beta(\gamma)(e_\gamma+e_{-\gamma})$ (GRH, NCZ; \file{density2/PRIMECORR.md}, Lemma~P0$'$).

\begin{hypothesis}[H$_{\rm dec}(L;s_1,s_2,g,T^*)$, prime side]
\textup{(Z)} No zero has ordinate in $(s_1-g,s_1+g)\cup(s_2-g,s_2+g)$. \textup{(N)} The band contains a zero: $n_I:=\sum_{s_1<\gamma<s_2}m_\gamma\ge1$. \textup{(D)} For all $a\ge0$, $b\ge1$ with $a+b\le L$, the log-averaged joint moment $M_{a,b}=\lim_Y\frac1Y\int_1^Y\Pi_R^a\Pi_I^b\,dy$ satisfies $(-1)^{a+b}M_{a,b}-(-1)^aM_{a,0}\,\mu_b(n_I)<1$, where $\mu_b(n)=E[(\sum_{i\le n}2\cos\theta_i)^b]$.
\end{hypothesis}

Each quantity in H$_{\rm dec}$ is a limit of prime sums: (Z) and (N) are one-level statements via the explicit formula. Under (Z), the excess in (D) is a sum of nonnegative weights over sequences whose net vector is a relation involving the band, each relation among zeros of height $\le T^*$ having weight exactly $1$ (\file{density2/PRIMECORR.md}, Lemma~P1$'$); in particular (D) at $(a,b)=(0,2)$ forces the zeros in the band to be simple.

\begin{theorem}[GRH, NCZ, band decoupling; \rd{2 reads}]\label{thm:P7}
Suppose that along a sequence $L_j\to\infty$ the hypothesis H$_{\rm dec}(L_j;s_{1,j},s_{2,j},g_j,T^*_j)$ holds, including \textup{(N)}, with $T^*_j=\lfloor L_j/2\rfloor^4$, $2s_{1,j}\le s_{2,j}\le3s_{1,j}$ and $s_{1,j}/\log(s_{1,j}+3)=o(L_j)$. Then $\nu_M$ has no atoms, for every $M>0$.
\end{theorem}

\begin{proof}[Proof (outline; \file{density2/PRIMECORR.md}, \S3)]
If $s_{1,j}$ stays bounded along a subsequence, some zero $\gamma_0$ lies in infinitely many bands, and (D) with $L_j,T^*_j\to\infty$ excludes every relation involving $\gamma_0$; so $\gamma_0$ is simple and $\gamma_0\notin\spanQ{\Gamma\setminus\{\gamma_0\}}$, TS holds, and Theorem~\ref{thm:TS} applies. If $s_{1,j}\to\infty$: by Lemma~P1$'$ and (D) no relation of length $\le L$ among zeros of height $\le T^*$ involves the band, so the band, which contains $\gg s_1$ distinct ordinates, is free; Proposition~\ref{prop:Bprime} with the density bound $D\le C(1+(s_1/\log s_1)^{1/2})$ and the truncation at $T^*=N^4$ give $\nu_M(\{c\})\le C(1+(s_1/\log s_1)^{1/2})(\log^2L/L+L^{-1/2})+o(1)\to0$.
\end{proof}

The first read of this theorem (\file{density2/REFEREE-4b.md}) found that its first version, without (N), held vacuously for an empty band and was therefore false as stated; condition (N) and the treatment of bounded bands are the corrections it proposed. A second, independent read (\file{density2/REFEREE-5.md}) found the corrected theorem, in both cases, correct as stated. LI with simple zeros implies H$_{\rm dec}$ for bands at every large height. For spectra with simple zeros (in one admissible band per scale), GLI implies H$_{\rm dec}$, so Theorem~\ref{thm:GLI} restricted to such spectra is a special case of Theorem~\ref{thm:P7}; without simplicity it is not (a $\QQ$-independent spectrum with all multiplicities $2$ satisfies GLI but fails (D)). Devin's Theorem~5.4(1) does not imply H$_{\rm dec}$ (\file{density2/PRIMECORR.md}, \S3). (The example offered in the working note to show that H$_{\rm dec}$ is strictly weaker than GLI, and that it does not imply Devin's Theorem~5.4(1), fails condition (D), as the second read found; we make neither claim.)

\begin{remark}[heuristic; not a theorem]
Pointwise Hardy--Littlewood asymptotics for $k$-tuples of shifted primes with power-saving error terms, as used by Montgomery and Soundararajan \cite{MS}, do not reach these moments: for windows of fixed relative width $1/T$ the error terms exceed the main term by a factor $\asymp(x/T)^{k/2}x^{-1/2+\eps}$. And by Corollary~\ref{cor:GLIatom} the relations an atom needs may all be far longer than any fixed moment order, so a hypothesis on moments of fixed order $K$ (or $K=o(T/\log T)$), with any precision, cannot exclude them.
\end{remark}

\section{Further facts under GRH alone, and a reduction to finite models}\label{sec:reduction}

\subsection{Smoothing, level sets and bands}

\begin{theorem}[analytic smoothing; GRH, NCZ; \rd{1 read}]\label{thm:smooth}
Let $k:\RR\to\RR$ be measurable with $\int(1+|u|)|k(u)|du<\infty$, $|\hat k(\xi)|\le Ce^{-a|\xi|}$ for some $a>0$, and $\hat k(\gamma_0)\ne0$ for some $\gamma_0\in\Gamma$. Then $E^k_M(y)=\int E_M(e^{y-u})k(u)\,du$ (with $E_M(e^v)=0$ for $v<\log2$) has a limiting distribution, the law of the continuous function $F^k_M=\hat k(0)+\re\sum_\gamma m_\gamma r_\gamma\hat k(\gamma)e_\gamma$, and this law is absolutely continuous.
\end{theorem}

\begin{proof}
The limiting distribution follows from the $B^2$-approximation of $E_M(e^y)$ by $F_{M,T}(\iota(y))$ \cite[Theorem~1.3(a)]{GO}, which survives convolution with $k$ because $\int(1+|u|)|k|<\infty$, and from (I1) for the continuous function $F^k_M$. For every $\theta$, $t\mapsto F^k_M(\tau_t\theta)$ extends holomorphically to $|\im t|<a$ and is non-constant (its Bohr coefficient at $\gamma_0$ is $m_{\gamma_0}\frac{2M}{M+i\gamma_0}\hat k(\gamma_0)\theta_{\gamma_0}\ne0$), so Lemma~\ref{lem:analytic}, invariance of $\hA$ under $\tau_t$ and Fubini give absolute continuity.
\end{proof}

So every analytic mollification of the race in $\log x$ has a density under GRH alone; the obstacle is the lack of smoothness at scale $0$.

\begin{proposition}[no level set has interior; GRH, NCZ; \rd{1 read}]\label{prop:interior}
For every $M>0$, $c\in\RR$ and every nonempty open $U\subset\cA$: $\hA(U\setminus\{F_M=c\})>0$.
\end{proposition}

\begin{proof}
Suppose $F_M=c$ a.e.\ on $U$, and choose nonempty open $U''\subset U'\subset U$ and $\delta,\delta_0>0$ with $\tau_tU'\subset U$ for $|t|<\delta$ and $\tau_tU''\subset U'$ for $|t|<\delta_0$. For $\eta\in C_c^\infty((-\delta,\delta))$ put $W_\eta(\theta)=\frac12\int\eta+\sum_{g\in\pm\Gamma}m_g\theta_g\int\eta(t)e^{igt}dt$, a continuous function. Comparing coefficients, $\int F_M(\tau_t\theta)(-\eta'(t)+M\eta(t))dt=2MW_\eta(\theta)$ a.e., and the left side is $cM\int\eta$ for a.e.\ $\theta\in U'$; so $2W_\eta=c\int\eta$ on $U'$ (by continuity). By (I1) there are arbitrarily large $y$ with $\iota(y)\in U''$, hence a $t_0=\log p$ ($p\nmid q$ prime) with $\iota(t_0)\in U'$. At $\theta=\iota(t_0)$, (I5) with $\eta(\cdot-t_0)$ gives $\frac c2\int\eta=\frac12\int\eta-\sum_n\Lambda(n)\chi(n)n^{-1/2}\eta(\log n-t_0)+\int\eta(t-t_0)\Omega(t)dt$. Take $\eta(t)=\eta_1(t/\eps)$ with $\eta_1\ge0$, $\eta_1(0)=1$, and let $\eps\to0$: every term tends to $0$ except $-\chi(p)(\log p)/\sqrt p\ne0$.
\end{proof}

\begin{proposition}[finite-order smoothing creates no atoms; GRH, NCZ; \rd{1 read}]\label{prop:finiteorder}
Let $\kappa>0$ and $K_\kappa f(\theta)=\int_0^\infty\kappa e^{-\kappa u}f(\tau_{-u}\theta)du$, and take the continuous version of $K_\kappa F_M$. Then $\{K_\kappa F_M=c\}\subset\{F_M=c\}$ up to a null set, for every $c$.
\end{proposition}

\begin{proof}
For a.e.\ $\theta$, $g_\theta(t)=K_\kappa F_M(\tau_t\theta)$ is locally absolutely continuous with $g_\theta'=\kappa(f_\theta-g_\theta)$ a.e., where $f_\theta(t)=F_M(\tau_t\theta)$. The derivative of an absolutely continuous function vanishes a.e.\ on each level set, so $f_\theta=g_\theta=c$ for a.e.\ $t$ with $g_\theta(t)=c$; conclude by Fubini.
\end{proof}

\begin{proposition}[band densities; GRH, NCZ; \rd{1 read}]\label{prop:band}
$\nu_M(\{c\})=\lim_{\eps\downarrow0}\dlow\{x:|E_M(x)-c|<\eps\}=\lim_{\eps\downarrow0}\dup\{x:|E_M(x)-c|\le\eps\}$. So $\nu_{1/2}(\{0\})=0$ for $\chi_{-4}$ iff for every $\eta>0$ there is $\eps>0$ such that $\{x:|\pi(x;4,3)-\pi(x;4,1)|\le\eps\sqrt x/\log x\}$ has upper logarithmic density $<\eta$.
\end{proposition}

\begin{proof}
\eqref{eq:portmanteau}, and $\nu_M((c-\eps,c+\eps))\downarrow\nu_M(\{c\})$, $\nu_M([c-\eps,c+\eps])\downarrow\nu_M(\{c\})$.
\end{proof}

\begin{corollary}[ties; GRH; \rd{2 reads}]\label{cor:ties}
Assume GRH$(\chi_{-4})$. The set of exact ties $\pi(x;4,3)=\pi(x;4,1)$ has upper logarithmic density $\le\nu_{1/2}(\{0\})<0.0891$; under GRH together with any of the hypotheses of Section~\ref{sec:sufficient} it has logarithmic density $0$.
\end{corollary}

\begin{proof}
Ties lie in every band $\{|E_{1/2}|\le\eps\}$; apply \eqref{eq:portmanteau}, Theorem~\ref{thm:cert}, and the absence of atoms.
\end{proof}

\begin{remark}[ties]\label{rem:ties}
Problem~13 of \cite{HKMN} asks for an \emph{unconditional} proof that ties have density $0$; we do not address the unconditional question, and we do not prove density $0$ for ties, or for near-ties $|\pi(x;4,3)-\pi(x;4,1)|\le K$, under GRH alone. Excluding ties beyond the band would need a local limit statement for sums $\sum_{t<p\le t+u}\chi(p)$ over short intervals (\file{density/DENSITY.md}, Remark~6.4).
\end{remark}

Under GRH, an atom also forces balance in multiplicative short intervals (\file{density/DENSITY.md}, Proposition~6.2, \rd{1 read}): if $\nu_{1/2}(\{0\})=\eta>0$ for $\chi_{-4}$, then for every $\eps>0$ and small $h$ the inequality $|\theta(xe^h;4,3)-\theta(x;4,3)-\theta(xe^h;4,1)+\theta(x;4,1)|<2\eps e^{h/2}\sqrt x$ holds on a set of lower logarithmic density $\ge\eta-\omega(h)$, with $\omega(h)\to0$; this is at a scale $\eps$ arbitrarily small compared with the root mean square $\asymp(h\log(1/h))^{1/2}$ of the same quantity.

\subsection{Reduction to finite integer models}

\begin{lemma}[reduction to finite models; GRH, NCZ; \rd{1 read}]\label{lem:reduction}
$F_{M,T}$ is a trigonometric polynomial on the finite-dimensional torus $\cA_T$, the closure of $\{(e^{i\gamma y})_{\gamma\le T}\}$, and its law under $\lambda_{\cA_T}$ is that of $F_{M,T}\circ\pi_T$ under $\hA$. For all $\eps>0$ and $T$,
\[
\nu_M((-\eps,\eps))\le\lambda_{\cA_T}\{|F_{M,T}|<2\eps\}+V_M^{>T}/\eps^2 .
\]
\end{lemma}

\begin{proof}
$\{|F_M|<\eps\}\subset\{|F_{M,T}\circ\pi_T|<2\eps\}\cup\{|F_M-F_{M,T}\circ\pi_T|\ge\eps\}$, and Chebyshev's inequality.
\end{proof}

Let $\cC(M,A_0,s)$ be the class of finite \emph{integer models}: a finite $\Gamma'\subset(0,\infty)$, integers $m_\gamma\ge1$ with $\sum_{\gamma\in\Gamma'\cap(t,t+1]}m_\gamma\le A_0\log(t+3)$ for all $t\ge0$ and $\sum m_\gamma8M^2/(M^2+\gamma^2)\le s$, with $F_{\Gamma'}=1+\re\sum_{\gamma\in\Gamma'}m_\gamma r_\gamma(M)e_\gamma$ on its closure torus. No condition is imposed on the relations among the $\gamma$. Put $\omega(\eps)=\sup_{\cC(M,A_0,s)}\lambda\{|F_{\Gamma'}|<\eps\}$.

\begin{conjecture}[U; open]\label{conj:U}
For $M=\frac12$, some admissible constant $A_0$ in (I7) for $\chi_{-4}$, and $s=0.1556\ge\sigma_{1/2}^2(\chi_{-4})=0.155568\ldots$: $\omega(\eps)\to0$ as $\eps\to0$.
\end{conjecture}

By Lemma~\ref{lem:reduction}, Conjecture~\ref{conj:U} implies that under GRH$(\chi_{-4})$ alone $\nu_{1/2}(\{0\})=0$ and the logarithmic density of $\{x:\pi(x;4,3)>\pi(x;4,1)\}$ exists: every truncation $F_{1/2,T}$ lies in the class, whatever the relations among the true ordinates. In this form the unknown group generated by the zeros disappears, and what remains is a Littlewood--Offord-type anti-concentration statement with \emph{dependent} phases (Haar measure on an arbitrary subtorus). Two remarks on its sharpness (made by the referee of this lemma). Without integrality it fails if $A_0\gtrsim2$: Example~\ref{ex:real} below at $M=0.05$, rescaled to the classical weight on the spectrum $10\NN$, has $\sigma^2=0.1439$, positive real weights $\{4.868,\frac12\}$, an atom of mass $\frac13$, and obeys the counting law with constant about $1.96$. With multiplicities growing like $(\log\gamma)^{5/2}$ (and $A_0\log$ replaced by $A(\log)^{5/2}$) it fails by Theorem~\ref{thm:intatom} (that $\sigma^2\lesssim0.125$ in this rescaled example is the referee's estimate, not separately certified).

\begin{lemma}[compactness; no hypothesis; \rd{1 read}]\label{lem:compact}
Fix $M>0$, a compact set $\mathcal K$ of real sequences (product topology) with $\sup_{m\in\mathcal K}\sum_{n>N}m_n^2/n^2\to0$ and closed under truncation, and a compact $C\subset\RR\setminus(1+2\ZZ)$. For the cyclic models $F_m=1+\re\sum_nm_n\frac{4M}{M+in}e^{int}$ on $\TT$, the following are equivalent: \textup{(a)} $\lambda\{F_m=c\}=0$ for all $m\in\mathcal K$, $c\in C$; \textup{(b)} $\sup_{m\in\mathcal K,c\in C}\lambda\{|F_m-c|<\eps\}\to0$ as $\eps\to0$.
\end{lemma}

(Examples: all integer sequences with $|m_n|\le B$, or with $0\le m_n\le A\log(n+2)$. Proof outline: a coordinatewise limit of near-atoms is an atom, by uniform $L^2$ tails and the portmanteau theorem; \file{density/DENSITY.md}, Lemma~9.18.)

\begin{corollary}[a special case of Conjecture~\ref{conj:U}; no hypothesis; \rd{1 read}]\label{cor:Ucase}
For $Y\ge1$ let $\cC_V(M,A_0,s;Y)$ be the set of models in $\cC(M,A_0,s)$ in which some ordinate $\gamma\le Y$ is a vertex of $\mathrm{conv}(\pm\Gamma')$, and $V^*(M,A_0)=8M^2A_0^2\sum_{t\ge0}\log^2(t+3)/(M^2+t^2)$. Then for $0<\eps<1$ and real $c$,
\[
\sup_{\cC_V(M,A_0,s;Y)}\lambda\{|F_{\Gamma'}-c|<\eps\}\le\frac{\frac12\log\frac{M^2+Y^2}{4M^2}+\frac12\log^+((1-c)^2+V^*(M,A_0))+\frac1{2e}}{\log(1/\eps)} .
\]
\end{corollary}

\begin{proof}
Lemmas~\ref{lem:JN} and \ref{lem:smallball} with $|m_\gamma|\ge1$ at the vertex, and $V_{\Gamma'}\le V^*$ (group the ordinates into unit intervals and use the counting law).
\end{proof}

The class $\cC_V$ is the ``easy'' part: it contains every model with one ordinate $\le Y$ outside the $\QQ$-span of the others, and no rank-one model with top ordinate $>Y$. In rank one with a dense spectrum the Jensen--Newton bound loses a full $\log T$, and this loss is real for integer models (\file{density3/MAHLER.md}, Propositions~5.1, 5.2); the vertex data used by Lemma~\ref{lem:JN} are the same in the real-weight class, where the analogue of Conjecture~\ref{conj:U} fails, so arguments based on them alone cannot prove it.

\section{Barriers}\label{sec:barriers}

Everything GRH gives about the hull function is its Fourier series: frequencies $\pm\Gamma$ and coefficients $m_\gamma\frac{2M}{M\pm i\gamma}$, i.e.\ $(D+M)F_M=2Mw$ with $D$ the derivative along the flow and $w=\frac12+\sum_{g\in\pm\Gamma}m_ge_g$, a formal series with \emph{positive integer} coefficients and constant term $\frac12$. The arguments of Sections~\ref{sec:jump} and \ref{sec:reversal} use only the positivity of the weights and the sizes of $\sigma_M^2$ and $V_M$ (Theorems~\ref{thm:nr}, \ref{thm:cert} use in addition finite non-resonance and certified zeros), and those of Section~\ref{sec:sufficient} use hypothetical arithmetic information on $\Gamma$. This section shows what the available structure cannot decide. All statements here are about model functions and use no hypothesis. The \emph{cyclic model} with real weights $m_n$ ($n\ge1$, $\sum m_n^2/n^2<\infty$) is
\[
F=1+\re\sum_{n\ge1}m_n\frac{4M}{M+in}e^{int}\quad\text{on }\TT=\RR/2\pi\ZZ ,
\]
the analogue of $F_M$ with the ordinates replaced by the positive integers. Rescaling $t\mapsto t/\Omega$, $\Omega=1/(2M)$, turns the model at weight $M$ on the frequencies $n$ into the model at the classical weight $\frac12$ on the frequencies $\Omega n$, with the same weights, the same $\sigma^2=\sum m_n8M^2/(M^2+n^2)$ and the same law.

\subsection{Real weights}

\begin{proposition}[\rd{1 read}]\label{prop:model71}
Let $M>0$, $0<a<\pi$ with $a/\pi\notin\QQ$, $\kappa^2=2\pi^2/(M^2a^2)$, and $m_n=\kappa^2(M^2+n^2)\sin^2(na/2)/(\pi^2n^2)>0$. Then $F(t)=\frac{M\kappa^2}\pi[M(a-|t|)+\operatorname{sgn}t]\mathbf 1_{|t|<a}$ on $(-\pi,\pi]$; so $F=0$ on the arc $a<|t|\le\pi$, the law of $F$ has an atom at $0$ of mass $1-a/\pi$, and $\sigma^2=4\pi/a-2>2$.
\end{proposition}

\begin{proof}
With $v=\kappa\mathbf 1_{[0,a]}$ and $\Phi(t)=\frac1{2\pi}\int v(t+s)v(s)ds=\kappa^2(a-|t|)^+/(2\pi)$, $\hat\Phi(n)=|\hat v(n)|^2$; the function $2M(M\Phi-\Phi')$ equals the displayed closed form, its coefficients are $2M(M-in)|\hat v(n)|^2=m_n\cdot2M/(M+in)$ for $n\ne0$ and $1$ for $n=0$.
\end{proof}

\begin{proposition}[\rd{1 read}]\label{prop:model72}
For every closed arc $J\subset(0,\pi)$ at positive distance from $0$ there is $M_1>0$ such that for every $0<M\le M_1$ there are positive weights $m_n>0$ ($n\ge1$) with $\sigma^2\le C'M$ and $F=0$ on $J$. Rescaled, the same weights on the frequencies $\Omega\NN$ give models at the classical weight $\frac12$ with positive weights and $\sigma^2$ arbitrarily small, with an atom at $0$.
\end{proposition}

(The proof, by solving $\Psi'=M\Psi$ on $J$ for $\Psi=1+\sum c_n\cos nt$ with a positive-definite dominating bump, is in \file{density/DENSITY.md}, Proposition~7.2. \emph{Remark, not re-read:} the version there also gives $1\le m_n\le C/M$ for every $n$; these extra inequalities were added after the read.)

\begin{proposition}[the classical weight; computer-assisted; \rd{1 read}]\label{prop:model73}
For $M=\frac12$ there are $c_n\ge0$ ($n\ge1$) with $\sum c_n<0.2717$ such that $F(t)=1+\sum_nc_n(\cos nt+2n\sin nt)$ vanishes on the arc $J=[2,2.3]$. Equivalently: with the bounded weights $m_n=c_n(\frac14+n^2)\ge0$, the cyclic model at $M=\frac12$ has $\sigma^2=2\sum c_n<0.5434<2$ and an atom at $0$ of mass $\ge0.3/2\pi>0.047$.
\end{proposition}

\begin{proof}[Proof (outline; \file{density/DENSITY.md}, Proposition~7.3)]
Start from the trigonometric polynomial $\Psi_N=1+\sum_{n\le100}c^{(N)}_n\cos nt$, with $33$ fitted positive rational coefficients and $F_N=\Psi_N-2\Psi_N'$; certified Taylor bounds give $\sup_J|F_N|\le2.16\cdot10^{-15}$ and $\sup_J|F_N'|\le4.91\cdot10^{-10}$. Solve $\Psi_c-\Psi_c'/M=-F_N$ exactly on $J$, extend $\Psi_c$ to a $C^1$ even function supported near $J$ with explicit bounds on its cosine coefficients, and add $\eps P$ with $P(t)=((1-|t|/r)^+)^2$, $r=1.9$, $\eps=1.4\cdot10^{-10}$, whose cosine coefficients dominate the correction's; after normalisation all coefficients are $\ge0$ and $\Psi-\Psi'/M=0$ on $J$. All comparisons are in Arb (\file{density/model_certified.py}); the read re-ran it with byte-identical output and checked the construction independently at $50$ digits.
\end{proof}

\begin{remark}[the model test]
Fix the weights of Proposition~\ref{prop:model73} and consider the whole family $F^{\rm mod}_M=1+\re\sum m_n\frac{4M}{M+in}e^{int}$, $M>0$. It has a compact connected hull with dense flow, positive bounded weights, the identities between weights of Section~\ref{sec:jump} (since $\sigma^2<2$ near $M=\frac12$), real-analyticity in $M$, a critical weight, time reversal and everything in Theorems~\ref{thm:reversal}, \ref{thm:cheb52}; and its law at $M=\frac12$ has an atom at $0$. So no argument using only these features can prove $\nu_{1/2}(\{0\})=0$. The model lacks: ($\alpha$) integer weights; ($\beta$) the explicit formula on a dense orbit (Proposition~\ref{prop:interior}; here the level set is open); ($\gamma$) weights $\ge1$ on every frequency (the unrefereed strengthening of Proposition~\ref{prop:model72} has them, on a sparse spectrum); ($\delta$) a non-cyclic group $G$ (the true $G$ is not cyclic, by the explicit formula); ($\epsilon$) the growth of $N(T)$, hence unbounded $F_M$. A proof must use at least one of these.
\end{remark}

\subsection{What integrality gives}\label{sec:integrality}
\emph{Nested combs.} Let $1=N_0\mid N_1\mid\cdots$ with $N_{k+1}\ge2N_k$, $b_k\in\RR$ with $\sum|b_k|/N_k<\infty$, $m_n=\sum_{k:N_k\mid n}b_k$ and $B(J)=\sum_{k\le J}b_k$. With $h_k=2\pi/N_k$ and the exponential sawtooth $\rho_k(t)=h_ke^{-Ms}/(1-e^{-Mh_k})-1/M$, $s=t\bmod h_k$, one has $(D+M)\rho_k=h_k\sum_j\delta_{jh_k}-1$ and $F=1+2M\sum_kb_k\rho_k$, uniformly convergent.

\begin{theorem}[parity rigidity; no hypothesis; \rd{1 read}]\label{thm:nested}
If all $b_k\in\ZZ$, then $\lambda\{F=c\}=0$ for every $c\notin1+2\ZZ$; in particular the law of $F$ has no atom at $0$. If moreover infinitely many $b_k\ne0$, the law of $F$ has no atoms at all.
\end{theorem}

\begin{proof}
Suppose $S=\{F=c\}$ has positive measure. By the martingale convergence theorem, for a.e.\ $t_*\in S$ the level-$J$ cell $I\ni t_*$ ($I=[ah_J,(a+1)h_J)$) satisfies $|S\cap I|\ge\frac78|I|$ for all large $J$. On the interior of $I$ no comb with $k\le J$ has an atom, so $F(t)=1-2B(J)+2M\alpha e^{-Mt}+2M\Phi(t)$ with $\alpha=\sum_{k\le J}b_kA_k$ a constant and $\Phi=\sum_{k>J}b_k\rho_k$, which is $h_{J+1}$-periodic with $|\Phi|\le\eps_J=\sum_{k>J}|b_k|h_k\to0$. As $I$ consists of $r\ge2$ sub-cells of length $h'=h_{J+1}$, there is $t$ with $t,t+h'\in S\cap I$; then $F(t)=F(t+h')$ and $\Phi(t)=\Phi(t+h')$ force $\alpha=0$. So $|c-1+2B(J)|\le2M\eps_J$ with $B(J)\in\ZZ$, for all large $J$; letting $J\to\infty$, $c\in1+2\ZZ$, and $B(J)$ is eventually constant, so $b_k=0$ for large $k$.
\end{proof}

\begin{example}[real weights; \rd{1 read}]\label{ex:real}
Take $N=(1,3)$, any $M>0$, $\mathfrak a_0=2\pi/(1-e^{-2\pi M})$, $\mathfrak a_1=\frac{2\pi}3e^{4\pi M/3}/(1-e^{-2\pi M/3})$, $b_0=\mathfrak a_1/(2(\mathfrak a_1-\mathfrak a_0))$, $b_1=\frac12-b_0$. Then $F=0$ on $(4\pi/3,2\pi)$, so the law has an atom $\ge\frac13$ at $0$; the weights $m_n\in\{b_0,\frac12\}$ are positive ($\mathfrak a_1/\mathfrak a_0=(1+e^v+e^{2v})/3>1$, $v=2\pi M/3$). At $M=\frac12$, $\sigma^2\approx1.842<2$.
\end{example}

So in one and the same class real nonnegative weights give an atom at $0$ and integer weights do not. The mechanism is the half-integer constant $\frac12$ of $w$ (which comes from the prime squares, i.e.\ from Chebyshev's bias itself) against the integrality of the partial sums $B(J)$.

\begin{theorem}[integrality excludes flat arcs; cyclic model; no hypothesis; \rd{1 read}]\label{thm:arcs}
Let $m_n\in\ZZ$ with $|m_n|\le C(1+n)^A$, let $F$ be the distribution with $\hat F(0)=1$, $\hat F(\pm n)=2Mm_n/(M\pm in)$, and suppose $F=c$ on a nonempty open arc $I$, as distributions. Then $g(z)=\sum_{n\ge1}m_nz^n$ is rational with poles at roots of unity in $\TT\setminus I$, and $c\in1+2\ZZ$. For rank-one hulls ($G\subset\QQ\gamma_0$) with integer $m_\gamma$ of polynomial growth and $\sum m_\gamma^2/\gamma^2<\infty$: if $F=c$ a.e.\ on a nonempty open set, then $c\in1+2\ZZ$.
\end{theorem}

\begin{proof}[Proof (outline; \file{density/DENSITY.md}, Theorem~9.16 and Corollary~9.16$'$)]
$w=(D+M)F/2M$ equals $c/2$ on $I$, so $\mathcal T=w-c/2$ is supported in $K=\TT\setminus I$. Its Herglotz transform $H(z)=\langle\mathcal T,\phi k_z\rangle$ ($k_z(t)=\frac{e^{it}+z}{e^{it}-z}$, $\phi$ a smooth cutoff equal to $1$ near $K$) is analytic on $\hat\CC\setminus K$, equals $2g(z)-\kappa$ for $|z|<1$ and $\kappa-2g(1/z)$ for $|z|>1$, $\kappa=(c-1)/2$. So $g$ continues across $I$, and by the P\'olya--Carlson theorem \cite{Carlson,Polya} it is rational; by Fatou's lemma on rational power series with integer coefficients $g=R/Q$ with $R,Q\in\ZZ[z]$, $Q(0)=1$, poles on $K\subset\TT$, so the leading coefficient of $Q$ is $\pm1$ and $\kappa=g(\infty)\in\ZZ$. The rank-one statement of the theorem follows by conditional expectation onto a circle factor.
\end{proof}

This is a folklore-level consequence of P\'olya--Carlson once $g$ is continued across the arc; we did not find the parity conclusion stated. For bounded real weights, continuation across an arc makes every right limit of the coefficients reflectionless there (Breuer--Simon \cite[Theorem~1.3]{BS}); no positive-measure version of these rigidity theorems is known, and the positive-measure version of the Duffin--Schaeffer theorem is stated as open by Borwein, Erd\'elyi and Littmann \cite{BEL}. With Theorem~\ref{thm:intatom} below it gives a sharp dichotomy for integer multiplicities: on arcs integrality alone forces parity, at any polynomial growth; on sets of positive measure with empty interior it does not. Related rigidity results, all proved:
\begin{itemize}
\item eventually periodic integer $m_n$ allow positive-measure level sets only at odd integers (\file{density/DENSITY.md}, Proposition~9.10; elementary, via Kronecker's theorem), so an atom at $0$ forces $g$ to have the circle as natural boundary; \rd{1 read}
\item the same for bounded integer $m_n$ that are ``Toeplitz-symmetric'' along infinitely many moduli $N$ (i.e.\ $n\mapsto m_n\mathbf 1[N\nmid n]$ is $N$-periodic on $\ZZ$) (\file{density/DENSITY.md}, Theorem~9.15), and for lacunary comb supports (Proposition~9.12 there); \rd{1 read each}
\item for one-sided Toeplitz $0/1$ (or integer) sequences with one hole per period at $0$, of any periods and patterns, a positive-measure level set forces symmetric patterns and an odd integer level; with bounded or eventually prime period ratios an atom forces $m$ periodic and level $\pm1$ (\file{density2/TOY.md}, Theorem~2.3, Corollary~2.4); \rd{2 reads, resp.\ 1}
\item if the driving distribution $w$ is a finite measure on any hull, integer weights allow positive-measure level sets only at odd integers (\file{density2/TOY.md}, Theorem~1.2, via Cohen's idempotent theorem \cite{Cohen}; cf.\ Helson \cite{Helson}); this never applies to the true race, whose $w$ is not a measure. \rd{2 reads}
\end{itemize}

\subsection{Integer multiplicities do not exclude atoms}\label{sec:intmodels}

\begin{theorem}[no hypothesis; \rd{1 read}]\label{thm:intatom}
For every $M>0$ there are integers $m_n\ge1$ with $m_n\ll_M(\log(n+2))^{5/2}$ such that the cyclic model $F$ vanishes a.e.\ on a set of measure $>0.1\cdot2\pi$.
\end{theorem}

\begin{proof}[Construction (outline; \file{density/DENSITY.md}, Theorem~9.6 and Lemma~9.7)]
Work in the class $F=(1-D/M)\Psi$, $\Psi$ even, in which $m_n=(M^2+n^2)\hat\Psi(n)/(2M^2)$ is real. Start from Example~\ref{ex:real} ($F_0=0$ on $J=(4\pi/3,2\pi)$). At scale $j$ ($N_j=4^j$, band $B_j=[N_j,N_{j+1})$) round the current multiplicities in $B_j$ \emph{up} to integers and cancel the rounding error $e^\Psi_j$ exactly by $\varphi_j=(1-D/M)(e^\Psi_j\beta_j/\hat\beta_j(0))$, where $\beta_j$ is a comb of bumps $b=c*\tilde c$ ($\hat b\ge0$) of period $2\pi/N'_j$, $N'_j=3N_{j+1}$, supported on a set $E_j$ of measure $2\pi\rho_j$, $\rho_j=0.05(j+1)^{-5/4}$. The comb has no Fourier coefficients in $(0,N'_j)$, so the multiplicities below $N_{j+1}$ become exactly integral and are never changed again; the corrections have nonnegative coefficients, so positivity is free; $\|\varphi_j\|_2\ll N_j^{-1/2}\rho_j^{-3/2}$ is summable; $F_\infty$ vanishes on $J\setminus\bigcup E_j$, of measure $>2\pi(\frac13-0.05\zeta(\frac54))>0.1\cdot2\pi$; and the decay of $\hat b$ gives $m_n\ll(\log n)^{5/2}$.
\end{proof}

The correction technique is of Menshov type (cf.\ Olevskii's survey \cite{Olevskii}); we know of no published statement of Theorem~\ref{thm:intatom} itself. The model violates the local zero-counting law $\sum_{\gamma\in[T,T+1]}m_\gamma\ll\log T$ (its local mass is $\asymp(\log T)^{5/2}$). Noise shaping removes this obstruction, at the price of an astronomical constant.

\begin{theorem}[noise-shaped integer atoms; no hypothesis; \rd{1 read}]\label{thm:noise}
For every $M>0$ and $u>0$ there are integers $m_n\ge1$ and a constant $C(M,u)$ with $m_n\le C(M,u)(\log(n+2))^u$ for all $n$, such that the cyclic model $F$ vanishes a.e.\ on a set $S$ of measure $\ge3\pi/(4q)$, $q=q(u)=\max(60,\lceil\frac{15\pi}{16}2^{4/u}\rceil)$. $S$ has empty interior (as it must, by Theorem~\ref{thm:arcs}).
\end{theorem}

\begin{corollary}[\rd{1 read}]\label{cor:noise}
Let $0<u\le1$ and $m_n$, $F$ as in Theorem~\ref{thm:noise}. \textup{(i)} $\sum_{n\in[T,T+1]}m_n\le A_c\log(T+3)$ for all $T\ge0$, $A_c=2C(M,u)$. \textup{(ii)} If $0<M\le M_0(u)=\min(\frac1{10},(1000\kappa_bq(u))^{-1/2})$, with $\kappa_b\approx2.7005$ a bump constant computed by quadrature, the construction can be run so that also $\sigma^2<\frac12$. \textup{(iii)} Then, rescaled, the model at the classical weight $\frac12$ on the frequencies $\Omega\NN$, $\Omega=1/(2M)$, has positive integer multiplicities on every frequency, obeys the local counting law with constant $A_c$, has $\sigma^2<\frac12$, and its law has an atom at $0$. \textup{(iv)} The constant is astronomical: $A_c\gtrsim10^{14}$ already at $q=60$, $M=\frac12$, and larger for $M\le M_0$; rescaling cannot bring it down to the constant $A_0$ of the true zeros.
\end{corollary}

\begin{proof}[Design (outline; complete proof in \file{density/DENSITY.md}, \S9.22)]
As in Theorem~\ref{thm:intatom}, but the rounding errors are shaped by a sigma--delta filter $(1-z)^{k_j}$ of slowly growing order (cf.\ \cite{DD,Gunturk}) and multiplied by a smooth cutoff $\varphi$ on $\TT$ equal to $1$ on the target arc and $0$ near the point $0$. A discrete Leibniz rule (Lemma~B there) bounds the shaped error on the arc by the $\ell^1$ norms of the Fourier coefficients of $(e^{it}-1)^r\varphi(t)$, which are $\le\tau^r\Lambda_r$ for explicit $\Lambda_r$ of polynomial growth, with $\tau<\frac1{10}$ ($|e^{it}-1|\le\tau$ on the support of $\varphi$); meanwhile multiplicities grow only like $2^{k_j}\ll(j+1)^{u'}$, $u'\le u$. A smooth slowly varying offset $\asymp(\log n)^{u'}$ keeps all multiplicities positive; its leak is removed with the same comb; the measure budget $\sum\lambda(E_j)$ is half of the arc.
\end{proof}

So integrality, positivity, the local counting law with a free constant, and $\sigma^2<2$ (in particular the refined question asked in \file{density/DENSITY.md} after Remark~9.9) do not exclude an atom at $0$. What is left on this line is the fixed-constant form (Conjecture~\ref{conj:U}) and bounded multiplicities, which GRH does not supply. The read of Theorem~\ref{thm:noise} re-derived all its lemmas, tested the key Lemma~B against worst-case errors (ratio $\le0.076$ of the bound), and ran a toy cascade with multiplicities integral to $3.6\cdot10^{-12}$, $\min m_n=1$, and vanishing on $S$ to $3.6\cdot10^{-17}$; it found gaps, all fixed, and nothing fatal.

\begin{proposition}[LI spectra are dense; no hypothesis; \rd{1 read}]\label{prop:LIdense}
Let $\Gamma\subset(0,\infty)$ be countable and discrete, $m:\Gamma\to\ZZ_{\ge1}$ and $\delta:\Gamma\to(0,\frac12)$. There is an injection $\gamma\mapsto\gamma'$ with $|\gamma'-\gamma|<\delta(\gamma)\min(1,\gamma/2)$ such that $\{\gamma'\}$ is $\QQ$-linearly independent. If $\sum m_\gamma^2/\gamma^2<\infty$, the model $1+\re\sum m_\gamma r_{\gamma'}(M)e_{\gamma'}$ (at least three terms) has an absolutely continuous law; the counting law $\sum_{(t,t+1]}m\le A\log(t+3)$ persists with $2A$ in place of $A$ (after adjusting the additive constant inside the logarithm), and $\sigma_M^2$ changes by a factor in $[(1+\delta_\infty)^{-2},(1-\delta_\infty)^{-2}]$, $\delta_\infty=\sup\delta$.
\end{proposition}

\begin{proof}
Choose $\gamma'_j$ successively outside the countable set $\spanQ{\gamma'_1,\dots,\gamma'_{j-1}}$. By Kronecker's theorem the hull is the full torus with product Haar measure; the characteristic function is $\prod J_0(\xi m_\gamma|r_{\gamma'}|)$, integrable as soon as three factors are kept.
\end{proof}

\begin{corollary}[barrier; \rd{1 read}]\label{cor:barrier}
The data ``positive integer multiplicities, the local counting law (with a free constant), $\sigma^2<2$'' do not decide whether the law has an atom at $0$: Corollary~\ref{cor:noise} gives a model with an atom, and Proposition~\ref{prop:LIdense} gives arbitrarily close spectra with the same data (constant $2A_c$) and absolutely continuous laws. So no property shared by the models of Theorem~\ref{thm:noise} and stable under small perturbations of the ordinates can imply ``no atom''.
\end{corollary}

This does not show that a proof must use unstable information (such as the explicit formula as an exact identity, which excludes cyclic spectra): a stable property of the true zeros that these lattice spectra violate, such as $N(t)=\frac t{2\pi}\log\frac{qt}{2\pi e}+O(\log t)$ with the true constant, might suffice.

\section{What is not proved}\label{sec:open}

\begin{itemize}
\item The existence of the logarithmic density of $\{x:\pi(x;4,3)>\pi(x;4,1)\}$ under GRH$(\chi_{-4})$ alone is \textbf{not} proved; nor $\nu_{1/2}(\{0\})=0$ for any explicitly given weight; nor the analogous statements for $\chi_{-3}$. The converse ``density exists $\Rightarrow$ no atom'' is not claimed.
\item Ties and near-ties $|\pi(x;4,3)-\pi(x;4,1)|\le K$ are not shown to have density $0$ under GRH alone (Remark~\ref{rem:ties}).
\item None of TS, FD, TS$_\ZZ$, V, GLI, H$_{\rm dec}$ is known for any $L$-function, and none can be verified by a finite computation. TS, FD and V each imply that two ordinates of $L(s,\chi)$ have an irrational ratio (the analogue of item~1 of \cite{HKMN}, open); TS$_\ZZ$ and GLI do not, as both hold for suitable sets of rational multiples of one number.
\item Conjecture~\ref{conj:U} is open, as is the case of bounded multiplicities (a relative of the open positive-measure Duffin--Schaeffer problem \cite{BEL}). Theorems~\ref{thm:intatom} and \ref{thm:noise} do not touch either: their multiplicities are unbounded and the counting constant is astronomical.
\end{itemize}

The state of the question, in one list. Under GRH alone: (1) $\nu_{1/2}(\{0\})=0$ is equivalent to the continuity at $m=0$ of the monotone function $m\mapsto\dlow\{D_m>0\}$, and it holds for all but countably many weights near $m=0$ (Theorem~\ref{thm:jump}); (2) the possible defect is $<0.0891$ ($q=4$), $<0.0653$ ($q=3$) (Theorem~\ref{thm:cert}); (3) an atom lives on a set with empty interior (Proposition~\ref{prop:interior}), invariant under every tame direction (Corollary~\ref{cor:tame}), and forces bounded-index integer relations with arbitrarily high ordinates (Corollary~\ref{cor:TSZatom}) and saturation of the spectrum by short relations (Corollary~\ref{cor:GLIatom}); (4) no argument using only the soft structure of the hull and the weights can exclude it (Proposition~\ref{prop:model73}), and integrality with a free counting constant cannot either (Corollary~\ref{cor:barrier}). A proof must use something these models lack. Candidates: the local counting law with the true constant (Conjecture~\ref{conj:U}); the true main term of $N(t)$ and the resulting unboundedness of $F_M$; a non-cyclic hull; the explicit formula on the dense orbit (which already excludes cyclic spectra) beyond Proposition~\ref{prop:interior}; or other arithmetic information on the zeros.

\section{Review, status and reproducibility}\label{sec:review}

All proofs, programs and literature searches were produced by AI (Anthropic's Claude models), and every result stated as proved passed at least one adversarial read by an independent AI agent, instructed to find errors and, where computations are involved, to use its own code. The reads, with the parts of the working notes they covered:
\begin{itemize}
\item \file{density/DENSITY.md} \S\S1--7: a first read (re-derived every statement; checked (I5) against stored zeros to $10^{-12}$; tested Theorem~\ref{thm:cheb52} on $232$ cyclic models and Theorem~\ref{thm:nr} on $206$ toy models by exact quadrature and $319$ adversarial models by Monte Carlo; re-ran \file{reflection_certified.py} with byte-identical output); a second read of the later additions (Propositions~\ref{prop:finiteorder}, \ref{prop:model72}, \ref{prop:model73}, the model test; re-ran \file{model_certified.py}); and an independent referee of Sections~\ref{sec:jump}, \ref{sec:reversal} (\file{density/REFEREE-2-DENSITY.md}: ``proof complete'' for both; all certified numbers reproduced independently). One false side statement was found and corrected (an invariance claim for arbitrary subsets of a level set, now stated for the full level set in Theorem~\ref{thm:tame}(b)), together with a gap in the range of the ``in particular'' bound of Theorem~\ref{thm:cheb52}, now closed.
\item \file{density/DENSITY.md} \S9 (integrality): five further reads (nested combs, Example~\ref{ex:real} and the conditional reduction to an integrality question, not used here; Theorem~\ref{thm:intatom}; Theorem~\ref{thm:arcs}, Proposition~9.17 there and Lemma~\ref{lem:compact}; a remark on integer certificates, found correct but of low use and not included here; Theorem~\ref{thm:noise}), and separate reads of Propositions~9.10 and 9.12, Lemma~9.13 (Lemma~\ref{lem:reduction} with its sharpness remarks) and Theorem~9.15 there, recorded as ``referee of this proposition'' in the note. None found a fatal error; the corrections are recorded in the note.
\item \file{density2/}: \file{HYP.md} (Section~\ref{sec:TSZ}) and \file{TOY.md} \S1 and Theorem~2.3 were read twice (\file{REFEREE-1-TOY.md}, \file{REFEREE-2.md}, \file{REFEREE-3.md}); \file{INVLO.md} and \file{PRIMECORR.md} twice (\file{REFEREE-4.md}, \file{REFEREE-4b.md}, and the independent \file{REFEREE-5.md}); the first read of Theorem~\ref{thm:P7} found its first version false for empty bands, and the second read found the corrected version, stated here, correct, while finding one comparison example in the note broken (not used here).
\item \file{density3/MAHLER.md} (Sections~\ref{sec:V}, \ref{sec:reduction}): one read (``repairable issues only''; all fixes applied).
\end{itemize}
Results of the working notes that did not pass a read are not included as results; two later strengthenings or rewrites are mentioned only as unrefereed remarks (the bound $m_n\ge1$ in Proposition~\ref{prop:model72}, and \file{density/DENSITY.md}, Example~3.3a). Heuristics are labelled as such.

\emph{Computations.} The certified bounds of Theorem~\ref{thm:cert} rest on: the zero lists of $L(s,\chi_{-4})$ and $L(s,\chi_{-3})$ certified complete and simple up to heights $640.1657$ and $700.0247$ (sign changes of Hardy's function and a winding-number count in Arb; \file{05-weighted-races-grh/cert/}); the certified variance bounds and the exact-integer non-resonance check of \cite{GO} (\file{05-weighted-races-grh/grhonly/}); and \file{density/reflection_certified.py} (Arb, $200$ bits). Proposition~\ref{prop:model73}: \file{density/model_certified.py}. Numerical checks of other statements are in the same directories, labelled as numerical.

\section*{Acknowledgements}
This paper was produced by AI (Anthropic's Claude models), directed by the author: the statements, proofs, programs, computations and literature searches were generated by AI. It was checked by independent AI referees as described in Section~\ref{sec:review}, and has not yet been refereed by a human expert. The companion manuscripts \cite{GO,MC,WR}, on which Sections~\ref{sec:setting}--\ref{sec:reversal} build, were produced in the same way. The author thanks the developers of FLINT/Arb and python-flint.


\begin{thebibliography}{99}\small

\bibitem{ANS} A. Akbary, N. Ng, M. Shahabi, \emph{Limiting distributions of the classical error terms of prime number theory}, Q. J. Math. \textbf{65} (2014), 743--780. arXiv:1306.1657.

\bibitem{AK} M. Aoki, S. Koyama, \emph{Chebyshev's bias against splitting and principal primes in global fields}, J. Number Theory \textbf{245} (2023), 233--262. arXiv:2203.12266.

\bibitem{Bailleul} A. Bailleul, \emph{Explicit Kronecker--Weyl theorems and applications to prime number races}, arXiv:2007.05763 (v3, 2022).

\bibitem{BEL} P. Borwein, T. Erd\'elyi, F. Littmann, \emph{Polynomials with coefficients from a finite set}, Trans. Amer. Math. Soc. \textbf{360} (2008), 5145--5154.

\bibitem{BS} J. Breuer, B. Simon, \emph{Natural boundaries and spectral theory}, Adv. Math. \textbf{226} (2011), 4902--4920. arXiv:1002.0823.

\bibitem{Cohen} P. J. Cohen, \emph{On a conjecture of Littlewood and idempotent measures}, Amer. J. Math. \textbf{82} (1960), 191--212.

\bibitem{Carlson} F. Carlson, \emph{\"Uber Potenzreihen mit ganzzahligen Koeffizienten}, Math. Z. \textbf{9} (1921), 1--13.

\bibitem{DD} I. Daubechies, R. DeVore, \emph{Approximating a bandlimited function using very coarsely quantized data: a family of stable sigma-delta modulators of arbitrary order}, Ann. of Math. (2) \textbf{158} (2003), 679--710.

\bibitem{Davenport} H. Davenport, \emph{Multiplicative Number Theory}, 3rd ed., Grad. Texts in Math. \textbf{74}, Springer, 2000.

\bibitem{Devin} L. Devin, \emph{Chebyshev's bias for analytic $L$-functions}, Math. Proc. Cambridge Philos. Soc. \textbf{169} (2020), 103--140. arXiv:1706.06394.

\bibitem{Devin2} L. Devin, \emph{Limiting properties of the distribution of primes in an arbitrarily large number of residue classes}, arXiv:1909.03975 (2019; journal version not checked).

\bibitem{Fiorilli} D. Fiorilli, \emph{Highly biased prime number races}, Algebra Number Theory \textbf{8} (2014), 1733--1767. arXiv:1210.6946.

\bibitem{Gunturk} C. S. G\"unt\"urk, \emph{One-bit sigma-delta quantization with exponential accuracy}, Comm. Pure Appl. Math. \textbf{56} (2003), 1608--1630.

\bibitem{HKMN} A. Hamieh, H. Kadiri, G. Martin, N. Ng, \emph{Comparative prime number theory problem list}, arXiv:2407.03530 (2024).

\bibitem{Hayani} M. Hayani, \emph{Chebyshev's bias without linear independence}, arXiv:2512.23302 (v3, 2026).

\bibitem{Helson} H. Helson, \emph{Note on harmonic functions}, Proc. Amer. Math. Soc. \textbf{4} (1953), 686--691.

\bibitem{HL} H. Helson, D. Lowdenslager, \emph{Prediction theory and Fourier series in several variables}, Acta Math. \textbf{99} (1958), 165--202.

\bibitem{IK} H. Iwaniec, E. Kowalski, \emph{Analytic Number Theory}, Amer. Math. Soc. Colloq. Publ. \textbf{53}, AMS, Providence, 2004.

\bibitem{Arb} F. Johansson, \emph{Arb: efficient arbitrary-precision midpoint-radius interval arithmetic}, IEEE Trans. Comput. \textbf{66} (2017), 1281--1292.

\bibitem{KR} J. Kaczorowski, O. Ramar\'e, \emph{Almost periodicity of some error terms in prime number theory}, Acta Arith. \textbf{106} (2003), 277--297.

\bibitem{MN} G. Martin, N. Ng, \emph{Inclusive prime number races}, Trans. Amer. Math. Soc. \textbf{373} (2020), 3561--3607. arXiv:1710.00088.

\bibitem{MS} H. L. Montgomery, K. Soundararajan, \emph{Primes in short intervals}, Comm. Math. Phys. \textbf{252} (2004), 589--617. arXiv:math/0409258.

\bibitem{Nazarov} F. L. Nazarov, \emph{Local estimates for exponential polynomials and their applications to inequalities of the uncertainty principle type}, Algebra i Analiz \textbf{5} (1993), no.~4, 3--66; English transl. St. Petersburg Math. J. \textbf{5} (1994), 663--717.

\bibitem{NV} H. Nguyen, V. Vu, \emph{Optimal inverse Littlewood--Offord theorems}, Adv. Math. \textbf{226} (2011), 5298--5319.

\bibitem{Olevskii} A. M. Olevskii, survey of modifications of functions and Fourier series (Menshov-type correction), Russian Math. Surveys (1985); title and data not re-checked.

\bibitem{Platt} D. Platt, \emph{Numerical computations concerning the GRH}, Math. Comp. \textbf{85} (2016). arXiv:1305.3087.

\bibitem{Polya} G. P\'olya, \emph{\"Uber gewisse notwendige Determinantenkriterien f\"ur die Fortsetzbarkeit einer Potenzreihe}, Math. Ann. \textbf{99} (1928), 687--706.

\bibitem{RS} M. Rubinstein, P. Sarnak, \emph{Chebyshev's bias}, Experiment. Math. \textbf{3} (1994), 173--197.

\bibitem{Rudin} W. Rudin, \emph{Fourier Analysis on Groups}, Interscience, New York, 1962.

\bibitem{Sheth} A. Sheth, \emph{Euler products at the centre and applications to Chebyshev's bias}, Math. Proc. Cambridge Philos. Soc. \textbf{179} (2025), 331--349. arXiv:2405.01512.

\bibitem{WR} R. Sharma, \emph{Rubinstein--Sarnak densities for weighted prime number races}, preprint, 2026, doi:10.5281/zenodo.21947101.

\bibitem{GO} R. Sharma, working note \file{05-weighted-races-grh/GRH-ONLY.md} in \cite{Repo} (all numbers cited as [GO, \dots] refer to this note); written up as \emph{Weighted prime races under GRH alone: lead inheritance, convex order and explicit bias bounds}, preprint, November 2026, \file{05-weighted-races-grh/grh-only-races.tex}, whose numbering differs.

\bibitem{MC} R. Sharma, \emph{Monotonicity and non-monotonicity of Chebyshev's bias in the weight}, preprint, November 2026; \file{05-weighted-races-grh/monotone-coupling.tex} in \cite{Repo}.

\bibitem{Repo} R. Sharma, \emph{Research-Math} (code, certificates, working notes and referee reports), \url{https://github.com/raghavs1729/Research-Math}, directory \texttt{2026-11-number-theory-q-series}.

\end{thebibliography}
\end{document}
